% Options for packages loaded elsewhere
\PassOptionsToPackage{unicode}{hyperref}
\PassOptionsToPackage{hyphens}{url}
%
\documentclass[
]{article}
\usepackage{amsmath,amssymb}
\usepackage{iftex}
\ifPDFTeX
  \usepackage[T1]{fontenc}
  \usepackage[utf8]{inputenc}
  \usepackage{textcomp} % provide euro and other symbols
\else % if luatex or xetex
  \usepackage{unicode-math} % this also loads fontspec
  \defaultfontfeatures{Scale=MatchLowercase}
  \defaultfontfeatures[\rmfamily]{Ligatures=TeX,Scale=1}
\fi
\usepackage{lmodern}
\ifPDFTeX\else
  % xetex/luatex font selection
\fi
% Use upquote if available, for straight quotes in verbatim environments
\IfFileExists{upquote.sty}{\usepackage{upquote}}{}
\IfFileExists{microtype.sty}{% use microtype if available
  \usepackage[]{microtype}
  \UseMicrotypeSet[protrusion]{basicmath} % disable protrusion for tt fonts
}{}
\makeatletter
\@ifundefined{KOMAClassName}{% if non-KOMA class
  \IfFileExists{parskip.sty}{%
    \usepackage{parskip}
  }{% else
    \setlength{\parindent}{0pt}
    \setlength{\parskip}{6pt plus 2pt minus 1pt}}
}{% if KOMA class
  \KOMAoptions{parskip=half}}
\makeatother
\usepackage{xcolor}
\usepackage[margin=1in]{geometry}
\usepackage{longtable,booktabs,array}
\usepackage{calc} % for calculating minipage widths
% Correct order of tables after \paragraph or \subparagraph
\usepackage{etoolbox}
\makeatletter
\patchcmd\longtable{\par}{\if@noskipsec\mbox{}\fi\par}{}{}
\makeatother
% Allow footnotes in longtable head/foot
\IfFileExists{footnotehyper.sty}{\usepackage{footnotehyper}}{\usepackage{footnote}}
\makesavenoteenv{longtable}
\usepackage{graphicx}
\makeatletter
\newsavebox\pandoc@box
\newcommand*\pandocbounded[1]{% scales image to fit in text height/width
  \sbox\pandoc@box{#1}%
  \Gscale@div\@tempa{\textheight}{\dimexpr\ht\pandoc@box+\dp\pandoc@box\relax}%
  \Gscale@div\@tempb{\linewidth}{\wd\pandoc@box}%
  \ifdim\@tempb\p@<\@tempa\p@\let\@tempa\@tempb\fi% select the smaller of both
  \ifdim\@tempa\p@<\p@\scalebox{\@tempa}{\usebox\pandoc@box}%
  \else\usebox{\pandoc@box}%
  \fi%
}
% Set default figure placement to htbp
\def\fps@figure{htbp}
\makeatother
\setlength{\emergencystretch}{3em} % prevent overfull lines
\providecommand{\tightlist}{%
  \setlength{\itemsep}{0pt}\setlength{\parskip}{0pt}}
\setcounter{secnumdepth}{-\maxdimen} % remove section numbering
\usepackage{graphicx}
\usepackage{booktabs}
\usepackage{float}
\usepackage{titling}

% Logo's bovenaan
\pretitle{%
  \begin{center}
  \includegraphics[width=4cm]{figures/BINF.png}\hfill
  \includegraphics[width=6cm]{figures/PAIH.png}\\[2cm]
  \huge
}
\posttitle{\end{center}}
\usepackage{booktabs}
\usepackage{longtable}
\usepackage{array}
\usepackage{multirow}
\usepackage{wrapfig}
\usepackage{float}
\usepackage{colortbl}
\usepackage{pdflscape}
\usepackage{tabu}
\usepackage{threeparttable}
\usepackage{threeparttablex}
\usepackage[normalem]{ulem}
\usepackage{makecell}
\usepackage{xcolor}
\usepackage{bookmark}
\IfFileExists{xurl.sty}{\usepackage{xurl}}{} % add URL line breaks if available
\urlstyle{same}
\hypersetup{
  pdftitle={In-depth Analysis for Lanaken},
  hidelinks,
  pdfcreator={LaTeX via pandoc}}

\title{In-depth Analysis for Lanaken}
\usepackage{etoolbox}
\makeatletter
\providecommand{\subtitle}[1]{% add subtitle to \maketitle
  \apptocmd{\@title}{\par {\large #1 \par}}{}{}
}
\makeatother
\subtitle{Leveraging Data Analytics for Efficient OR Planning at ZOL}
\author{}
\date{\vspace{-2.5em}October 2025}

\begin{document}
\maketitle

\begin{center}
\vspace{2cm}

Niels MARTIN\\
niels.martin@uhasselt.be\\[0.5em]
Maxim RIEBUS\\
maxim.riebus@uhasselt.be\\[0.5em]
Haroon THARWAT\\
haroon.tharwat@uhasselt.be\\

\vspace{6cm}

\textbf{Hasselt University} \\
Faculty of Business Economics \\
Research Group Business Informatics \\
Process Analytics in Healthcare
\end{center}

\newpage
\tableofcontents
\newpage

\section{Introduction}\label{introduction}

Efficient management of operating rooms (ORs) requires an understanding
of how planned surgical activity aligns with real-world execution. This
document provides an in-depth, site-specific analysis for \textbf{Campus
Lanaken}, serving as an analytical extension of the previous exploratory
study. The goal is to systematically address the key research questions
formulated in the project description:

\textbf{Research questions}

\begin{itemize}
\tightlist
\item
  \textbf{Planned versus observed timing}

  \begin{itemize}
  \tightlist
  \item
    How well do the planned operating durations correspond to the actual
    durations?\\
  \item
    What is the typical deviation between the planned start time and the
    actual entry into the operating room?\\
  \item
    Which specialisms or procedure types show the largest differences?\\
  \item
    How do these deviations scale relative to the planned duration
    (proportional or percentage-wise)?
  \end{itemize}
\item
  \textbf{Adherence to working hours}

  \begin{itemize}
  \tightlist
  \item
    How often do procedures run beyond regular working hours (08:00 to
    16:30)?\\
  \item
    What proportion of OR activity occurs in the evening or at night?\\
  \item
    How do these patterns differ between campuses?
  \end{itemize}
\item
  \textbf{Room allocation}

  \begin{itemize}
  \tightlist
  \item
    How often are procedures performed in a different room than
    originally planned?\\
  \item
    Are these reallocations concentrated in specific rooms or
    specialisms?
  \end{itemize}
\item
  \textbf{Urgency and scheduling}

  \begin{itemize}
  \tightlist
  \item
    How are urgent and non-elective cases distributed across campuses
    and time periods?\\
  \item
    To what extent do urgent surgeries overlap with elective activity
    and cause scheduling disruptions?
  \end{itemize}
\end{itemize}

Each of the following sections addresses one of these themes in depth,
focusing on the local characteristics of the Lanaken campus.

\subsection{Unique values and data
richness}\label{unique-values-and-data-richness}

The dataset captures a high degree of heterogeneity in surgical
practice. A total of 34889 unique patient IDs are recorded, spread
across 588 distinct procedure codes and 578 procedure names. Surgeon and
anesthesiologist coverage is broad, with 97 and 132 unique identifiers
respectively. These numbers reflect the diversity of clinical practice
at ZOL and underscore the importance of developing grouping strategies
to obtain actionable insights. Table \ref{tab:uniquevalues} provides an
overview of the number of unique values per variable.

\begin{table}[H]

\caption{\label{tab:unique values}Number of unique values per variable. \label{tab:uniquevalues}}
\centering
\begin{tabular}[t]{lr}
\toprule
Variable & UniqueValues\\
\midrule
PatientID & 34889\\
ProcedureCode & 588\\
ProcedureName & 578\\
AnesthesiologistID & 132\\
HeadSurgeonID & 97\\
\addlinespace
LengthOfStay & 66\\
PlannedOR & 9\\
ActualOR & 7\\
Weekday & 5\\
Year & 4\\
\addlinespace
AdmissionType & 3\\
\bottomrule
\end{tabular}
\end{table}

\subsection{Surgical procedure
durations}\label{surgical-procedure-durations}

Operating time is one of the most critical variables for planning.
Procedure durations display wide variation, both within and across
procedure types. Percentile analysis shows that 95\% of all procedures
finish within 86 minutes, 98\% within 114 minutes, 99\% within 133
minutes, and 99.9\% within 196 minutes. While the majority of procedures
are considerably shorter, the long right hand tail in the distribution
indicates the presence of outliers and highly complex cases that drive
variability in scheduling. Figure \ref{fig:duration} illustrates the
distribution of procedure durations, and Table
\ref{tab:durationpercentiles} summarizes the percentile values.

\begin{figure}[H]

{\centering \includegraphics[width=0.7\linewidth,]{In-Depth_Analysis_Lanaken_files/figure-latex/procedure_durations-1} 

}

\caption{Histogram of Surgery Duration \label{fig:duration}}\label{fig:procedure_durations}
\end{figure}

\newpage

\begin{longtable}[]{@{}lr@{}}
\caption{Percentiles of surgical procedure duration (in minutes).
\label{tab:durationpercentiles}}\tabularnewline
\toprule\noalign{}
Percentile & Duration (minutes) \\
\midrule\noalign{}
\endfirsthead
\toprule\noalign{}
Percentile & Duration (minutes) \\
\midrule\noalign{}
\endhead
\bottomrule\noalign{}
\endlastfoot
95\% & 86 \\
98\% & 114 \\
99\% & 133 \\
99.9\% & 196 \\
\end{longtable}

Looking at the most frequent procedures illustrates the large variation
in operating times. FACO EMULSIFICATIE has an average duration of 27.6
minutes and a median of 26 minutes, and transforaminele epidurale
infiltratie : lumbaal averages 12.6 minutes with a median of 12 minutes.
Facet radiofrequency lumbaal typically lasts 18.8 minutes on average
with a median of 18 minutes. In contrast, VITRECTOMIE have an average
duration of 55.3 minutes and a median of 53 minutes, highlighting the
substantial difference in resource requirements between routine short
procedures and major surgical interventions. This mix of short
standardized cases and longer complex operations underscores the
operational challenge of planning an efficient operating room schedule
that accommodates both types of activity. Table \ref{tab:top10dur}
summarizes the ten most frequent procedure types.

\begin{longtable}[]{@{}lrrr@{}}
\caption{Top 10 most frequent procedure types with mean and median
duration \label{tab:top10dur}}\tabularnewline
\toprule\noalign{}
ProcedureName & Count & Mean & Median \\
\midrule\noalign{}
\endfirsthead
\toprule\noalign{}
ProcedureName & Count & Mean & Median \\
\midrule\noalign{}
\endhead
\bottomrule\noalign{}
\endlastfoot
FACO EMULSIFICATIE & 11025 & 27.6 & 26 \\
transforaminele epidurale infiltratie : lumbaal & 7186 & 12.6 & 12 \\
facet radiofrequency lumbaal & 5577 & 18.8 & 18 \\
facet diagnostisch lumbaal & 4962 & 13.7 & 13 \\
proefsacrale infiltratie & 3290 & 11.5 & 11 \\
transforaminele epidurale infiltratie : sacraal & 2363 & 12.3 & 11 \\
epidurale infiltratie : cervicaal & 2091 & 14.0 & 13 \\
VITRECTOMIE & 1788 & 55.3 & 53 \\
facet radiofrequency cervicaal & 1693 & 22.8 & 22 \\
facet diagnostisch cervicaal & 1555 & 16.6 & 16 \\
\end{longtable}

\subsection{Admission type and
urgency}\label{admission-type-and-urgency}

The majority of procedures in the dataset are performed in a planned
setting. Ambulatory surgery accounts for 98.1\%, and inpatient
hospitalization for 1.9\%, while emergency admissions represent only one
recorded case. The urgency classification shows a similar pattern:
99.6\% of procedures are elective, and 0.4\% are non elective. Although
relatively small in number, the latter category remains important from a
planning perspective, as urgent procedures may disrupt elective
schedules and lead to cancellations or overtime. Table
\ref{tab:admissiontype} and Table \ref{tab:urgencytype} summarize both
distributions.

\begin{longtable}[]{@{}lrr@{}}
\caption{Frequency table for Admission Type
\label{tab:admissiontype}}\tabularnewline
\toprule\noalign{}
AdmissionType & Count & Percentage \\
\midrule\noalign{}
\endfirsthead
\toprule\noalign{}
AdmissionType & Count & Percentage \\
\midrule\noalign{}
\endhead
\bottomrule\noalign{}
\endlastfoot
DAG & 68108 & 98.1 \\
HOS & 1286 & 1.9 \\
SPOED & 1 & 0.0 \\
\end{longtable}

\begin{longtable}[]{@{}lrr@{}}
\caption{Frequency table for Urgency Classification
\label{tab:urgencytype}}\tabularnewline
\toprule\noalign{}
UrgencyType & Count & Percentage \\
\midrule\noalign{}
\endfirsthead
\toprule\noalign{}
UrgencyType & Count & Percentage \\
\midrule\noalign{}
\endhead
\bottomrule\noalign{}
\endlastfoot
electief & 69109 & 99.6 \\
niet-electief & 286 & 0.4 \\
\end{longtable}

\subsection{Temporal distribution of
activity}\label{temporal-distribution-of-activity}

Surgical activity is spread consistently across weekdays. Fridays
account for 18.5 percent of weekly activity, Tuesdays for 21.5 percent,
Thursdays for 19.1 percent, Wednesdays each for 19.8 percent, and
Mondays each for 21.2 percent. There are no Weekend days showen in the
data. Annual trends reveal steady growth, with the number of procedures
increasing from 20016 in 2022 to 20636 in 2024. At the time of data
extraction, 8702 procedures in 2025 represented 12.5 percent of the
total. Together, Table \ref{tab:weekday} and Table \ref{tab:year}
summarize the distribution of surgical activity across weekdays and
years.

\begin{longtable}[]{@{}
  >{\centering\arraybackslash}p{(\linewidth - 4\tabcolsep) * \real{0.1389}}
  >{\centering\arraybackslash}p{(\linewidth - 4\tabcolsep) * \real{0.1111}}
  >{\centering\arraybackslash}p{(\linewidth - 4\tabcolsep) * \real{0.1806}}@{}}
\caption{Frequency table for Day of the Week (sorted by descending
count) \label{tab:weekday}}\tabularnewline
\toprule\noalign{}
\begin{minipage}[b]{\linewidth}\centering
Weekday
\end{minipage} & \begin{minipage}[b]{\linewidth}\centering
Count
\end{minipage} & \begin{minipage}[b]{\linewidth}\centering
Percentage
\end{minipage} \\
\midrule\noalign{}
\endfirsthead
\toprule\noalign{}
\begin{minipage}[b]{\linewidth}\centering
Weekday
\end{minipage} & \begin{minipage}[b]{\linewidth}\centering
Count
\end{minipage} & \begin{minipage}[b]{\linewidth}\centering
Percentage
\end{minipage} \\
\midrule\noalign{}
\endhead
\bottomrule\noalign{}
\endlastfoot
Di & 14902 & 21.5 \\
Ma & 14710 & 21.2 \\
Wo & 13742 & 19.8 \\
Do & 13233 & 19.1 \\
Vr & 12808 & 18.5 \\
\end{longtable}

\begin{longtable}[]{@{}
  >{\centering\arraybackslash}p{(\linewidth - 4\tabcolsep) * \real{0.0972}}
  >{\centering\arraybackslash}p{(\linewidth - 4\tabcolsep) * \real{0.1111}}
  >{\centering\arraybackslash}p{(\linewidth - 4\tabcolsep) * \real{0.1806}}@{}}
\caption{Frequency table for Year \label{tab:year}}\tabularnewline
\toprule\noalign{}
\begin{minipage}[b]{\linewidth}\centering
Year
\end{minipage} & \begin{minipage}[b]{\linewidth}\centering
Count
\end{minipage} & \begin{minipage}[b]{\linewidth}\centering
Percentage
\end{minipage} \\
\midrule\noalign{}
\endfirsthead
\toprule\noalign{}
\begin{minipage}[b]{\linewidth}\centering
Year
\end{minipage} & \begin{minipage}[b]{\linewidth}\centering
Count
\end{minipage} & \begin{minipage}[b]{\linewidth}\centering
Percentage
\end{minipage} \\
\midrule\noalign{}
\endhead
\bottomrule\noalign{}
\endlastfoot
2022 & 20016 & 28.8 \\
2023 & 20041 & 28.9 \\
2024 & 20636 & 29.7 \\
2025 & 8702 & 12.5 \\
\end{longtable}

In summary, procedure durations show strong variability with a long
tailed distribution driven by complex surgeries. The majority of
procedures are elective, while non elective cases still have an
important impact on schedule stability. Activity is distributed evenly
across weekdays and has grown consistently in recent years. These
findings form the baseline against which more advanced performance
metrics such as start and end deviations, utilization rates, and
predictive models can be developed in subsequent sections of this
report.

\newpage

\section{How accurately does planning reflect reality in the
OR?}\label{how-accurately-does-planning-reflect-reality-in-the-or}

\subsection{Are planned and observed durations
aligned?}\label{are-planned-and-observed-durations-aligned}

The comparison between planned and observed surgery durations shows a
broad spread around the one to one line. Many procedures lie close to
the diagonal, indicating that their observed durations correspond
reasonably well to the planned values. At the same time, a substantial
number of points fall above or below this line, showing that both
underestimation and overestimation occur frequently. Most deviations
appear as longer than planned surgeries, visible in the dense area below
the red dashed line. The dispersion increases for longer procedures,
reflecting greater uncertainty for complex cases. Figure
\ref{fig:planned_vs_observed} displays this relationship.

\begin{figure}[h]

{\centering \includegraphics[width=0.7\linewidth,]{In-Depth_Analysis_Lanaken_files/figure-latex/planned vs observed-1} 

}

\caption{Planned vs. Observed Surgery Duration (Lanaken) \label{fig:planned_vs_observed}}\label{fig:planned vs observed}
\end{figure}

On average, surgeries at Lanaken show meaningful variation between
planned and actual durations. Among procedures that last longer than
planned, the mean deviation is 9 minutes with a median of 5 minutes,
while shorter cases deviate on average 7.1 minutes with a median of 5
minutes. The standard deviations of 11.2 minutes and 15.2 minutes
highlight the presence of both moderate deviations and substantial
outliers. While most differences remain within reasonable limits, a
small number of extreme values extend the overall range.

Almost half of all surgeries (41.7\%) take longer than planned, with an
interquartile range of 8 minutes, while 58.3\% finish earlier, with an
interquartile range of 7 minutes. This distribution shows that
deviations occur in both directions, but the spread is larger for cases
that finish earlier. The considerable variability within each category
indicates that although overall planning accuracy is reasonable at the
aggregate level, individual durations remain difficult to predict.
Improving the consistency of duration estimates, particularly for
procedures that frequently deviate, could help reduce daily schedule
variability and limit delays across the operating list. Table
\ref{tab:dur_diff_direction} provides a detailed overview of these
differences.

\begin{longtable}[]{@{}
  >{\centering\arraybackslash}p{(\linewidth - 20\tabcolsep) * \real{0.2421}}
  >{\centering\arraybackslash}p{(\linewidth - 20\tabcolsep) * \real{0.0737}}
  >{\centering\arraybackslash}p{(\linewidth - 20\tabcolsep) * \real{0.0947}}
  >{\centering\arraybackslash}p{(\linewidth - 20\tabcolsep) * \real{0.0737}}
  >{\centering\arraybackslash}p{(\linewidth - 20\tabcolsep) * \real{0.0632}}
  >{\centering\arraybackslash}p{(\linewidth - 20\tabcolsep) * \real{0.0632}}
  >{\centering\arraybackslash}p{(\linewidth - 20\tabcolsep) * \real{0.0632}}
  >{\centering\arraybackslash}p{(\linewidth - 20\tabcolsep) * \real{0.0632}}
  >{\centering\arraybackslash}p{(\linewidth - 20\tabcolsep) * \real{0.0737}}
  >{\centering\arraybackslash}p{(\linewidth - 20\tabcolsep) * \real{0.0842}}
  >{\centering\arraybackslash}p{(\linewidth - 20\tabcolsep) * \real{0.1053}}@{}}
\caption{Difference between actual and planned duration (minutes),
separated by direction \label{tab:dur_diff_direction}}\tabularnewline
\toprule\noalign{}
\begin{minipage}[b]{\linewidth}\centering
direction
\end{minipage} & \begin{minipage}[b]{\linewidth}\centering
Mean
\end{minipage} & \begin{minipage}[b]{\linewidth}\centering
Median
\end{minipage} & \begin{minipage}[b]{\linewidth}\centering
SD
\end{minipage} & \begin{minipage}[b]{\linewidth}\centering
IQR
\end{minipage} & \begin{minipage}[b]{\linewidth}\centering
Min
\end{minipage} & \begin{minipage}[b]{\linewidth}\centering
Max
\end{minipage} & \begin{minipage}[b]{\linewidth}\centering
P95
\end{minipage} & \begin{minipage}[b]{\linewidth}\centering
P99
\end{minipage} & \begin{minipage}[b]{\linewidth}\centering
P99\_9
\end{minipage} & \begin{minipage}[b]{\linewidth}\centering
Percent
\end{minipage} \\
\midrule\noalign{}
\endfirsthead
\toprule\noalign{}
\begin{minipage}[b]{\linewidth}\centering
direction
\end{minipage} & \begin{minipage}[b]{\linewidth}\centering
Mean
\end{minipage} & \begin{minipage}[b]{\linewidth}\centering
Median
\end{minipage} & \begin{minipage}[b]{\linewidth}\centering
SD
\end{minipage} & \begin{minipage}[b]{\linewidth}\centering
IQR
\end{minipage} & \begin{minipage}[b]{\linewidth}\centering
Min
\end{minipage} & \begin{minipage}[b]{\linewidth}\centering
Max
\end{minipage} & \begin{minipage}[b]{\linewidth}\centering
P95
\end{minipage} & \begin{minipage}[b]{\linewidth}\centering
P99
\end{minipage} & \begin{minipage}[b]{\linewidth}\centering
P99\_9
\end{minipage} & \begin{minipage}[b]{\linewidth}\centering
Percent
\end{minipage} \\
\midrule\noalign{}
\endhead
\bottomrule\noalign{}
\endlastfoot
Longer than planned & 9 & 5 & 11.2 & 8 & 1 & 331 & 29 & 53.5 & 108 &
41.7\% \\
Shorter than planned & 7.1 & 5 & 15.2 & 7 & 0 & 890 & 21 & 40 & 80.1 &
58.3\% \\
\end{longtable}

Expressing deviations as a percentage of the planned duration provides a
standardized view across short and long procedures. Most cases fall
within about ±40 percent of the planned time, indicating a generally
acceptable level of scheduling accuracy in routine practice. The
histogram peaks slightly below zero, showing that surgeries are on
average completed a bit earlier than planned. The shape is roughly
symmetrical, although the right tail extends further, indicating that
large overruns occur less frequently but can be substantial when they
appear. This pattern suggests that while most procedures are reasonably
well estimated, a smaller group of cases contributes disproportionately
to the overall variability in duration accuracy. Figure
\ref{fig:relative_dev} shows the full distribution.

\begin{figure}[h]

{\centering \includegraphics[width=0.7\linewidth,]{In-Depth_Analysis_Lanaken_files/figure-latex/relative_deviation_hist-1} 

}

\caption{Distribution of relative deviation between planned and observed duration \label{fig:relative_dev}}\label{fig:relative_deviation_hist}
\end{figure}

Duration deviations differ markedly between admission types.Inpatient,
HOS, cases show the highest variability, with an average deviation of
14.5 percent and a median of 6.7 percent, while the coefficient of
variation is 3.4 This means that even if the mean deviation seems
moderate, individual cases differ strongly from their planned duration,
which is typical for unpredictable acute interventions.

Inpatient, SPOED, surgery has a deviation of -17.8 percent and a median
of -17.8 percent, suggesting that this case was a rare occurance. The
coefficient of variation of cannot be calculated as there is only one
recorded case of an emrgency surgery.

Ambulatory, DAG, surgeries are the most predictable group. Their mean
deviation is 2 percent, the median is minus 5 percent, and the
coefficient of variation of 21.9 mainly reflects the fact that even
moderate absolute differences translate into larger percentages for
short procedures. These results show that standardized short-stay
procedures are generally estimated accurately in the planning process.
Table \ref{tab:rel_dev_admission} summarizes these patterns.

\begin{longtable}[]{@{}
  >{\centering\arraybackslash}p{(\linewidth - 10\tabcolsep) * \real{0.1538}}
  >{\centering\arraybackslash}p{(\linewidth - 10\tabcolsep) * \real{0.0769}}
  >{\centering\arraybackslash}p{(\linewidth - 10\tabcolsep) * \real{0.2885}}
  >{\centering\arraybackslash}p{(\linewidth - 10\tabcolsep) * \real{0.3077}}
  >{\centering\arraybackslash}p{(\linewidth - 10\tabcolsep) * \real{0.0865}}
  >{\centering\arraybackslash}p{(\linewidth - 10\tabcolsep) * \real{0.0865}}@{}}
\caption{Relative duration deviation and coefficient of variation by
admission type (Lanaken) \label{tab:rel_dev_admission}}\tabularnewline
\toprule\noalign{}
\begin{minipage}[b]{\linewidth}\centering
AdmissionType
\end{minipage} & \begin{minipage}[b]{\linewidth}\centering
n
\end{minipage} & \begin{minipage}[b]{\linewidth}\centering
Mean relative deviation (\%)
\end{minipage} & \begin{minipage}[b]{\linewidth}\centering
Median relative deviation (\%)
\end{minipage} & \begin{minipage}[b]{\linewidth}\centering
SD (\%)
\end{minipage} & \begin{minipage}[b]{\linewidth}\centering
CV (\%)
\end{minipage} \\
\midrule\noalign{}
\endfirsthead
\toprule\noalign{}
\begin{minipage}[b]{\linewidth}\centering
AdmissionType
\end{minipage} & \begin{minipage}[b]{\linewidth}\centering
n
\end{minipage} & \begin{minipage}[b]{\linewidth}\centering
Mean relative deviation (\%)
\end{minipage} & \begin{minipage}[b]{\linewidth}\centering
Median relative deviation (\%)
\end{minipage} & \begin{minipage}[b]{\linewidth}\centering
SD (\%)
\end{minipage} & \begin{minipage}[b]{\linewidth}\centering
CV (\%)
\end{minipage} \\
\midrule\noalign{}
\endhead
\bottomrule\noalign{}
\endlastfoot
HOS & 1286 & 14.5 & 6.7 & 49 & 3.4 \\
DAG & 68108 & 2 & -5 & 42.9 & 21.9 \\
SPOED & 1 & -17.8 & -17.8 & NA & NA \\
\end{longtable}

A similar pattern appears when comparing elective and non elective
surgeries. Non elective procedures have an average relative deviation of
-2.6 percent and a median of minus 7.3 percent, with a coefficient of
variation of -12.1, highlighting their unpredictable character. Elective
operations show smaller deviations, with an average of 2.2 percent and a
median of minus 2.2 percent. Their coefficient of variation of 19.5
results from the influence of a small number of extreme outliers rather
than consistent misestimation. Overall, these findings confirm that
while elective cases generally follow expected durations well, non
elective surgeries remain the main source of variability in the
operating schedule. Table \ref{tab:rel_dev_urgency} summarizes these
results.

\begin{longtable}[]{@{}
  >{\centering\arraybackslash}p{(\linewidth - 10\tabcolsep) * \real{0.1538}}
  >{\centering\arraybackslash}p{(\linewidth - 10\tabcolsep) * \real{0.0769}}
  >{\centering\arraybackslash}p{(\linewidth - 10\tabcolsep) * \real{0.2885}}
  >{\centering\arraybackslash}p{(\linewidth - 10\tabcolsep) * \real{0.3077}}
  >{\centering\arraybackslash}p{(\linewidth - 10\tabcolsep) * \real{0.0865}}
  >{\centering\arraybackslash}p{(\linewidth - 10\tabcolsep) * \real{0.0865}}@{}}
\caption{Relative duration deviation and coefficient of variation by
urgency type (Lanaken) \label{tab:rel_dev_urgency}}\tabularnewline
\toprule\noalign{}
\begin{minipage}[b]{\linewidth}\centering
UrgencyType
\end{minipage} & \begin{minipage}[b]{\linewidth}\centering
n
\end{minipage} & \begin{minipage}[b]{\linewidth}\centering
Mean relative deviation (\%)
\end{minipage} & \begin{minipage}[b]{\linewidth}\centering
Median relative deviation (\%)
\end{minipage} & \begin{minipage}[b]{\linewidth}\centering
SD (\%)
\end{minipage} & \begin{minipage}[b]{\linewidth}\centering
CV (\%)
\end{minipage} \\
\midrule\noalign{}
\endfirsthead
\toprule\noalign{}
\begin{minipage}[b]{\linewidth}\centering
UrgencyType
\end{minipage} & \begin{minipage}[b]{\linewidth}\centering
n
\end{minipage} & \begin{minipage}[b]{\linewidth}\centering
Mean relative deviation (\%)
\end{minipage} & \begin{minipage}[b]{\linewidth}\centering
Median relative deviation (\%)
\end{minipage} & \begin{minipage}[b]{\linewidth}\centering
SD (\%)
\end{minipage} & \begin{minipage}[b]{\linewidth}\centering
CV (\%)
\end{minipage} \\
\midrule\noalign{}
\endhead
\bottomrule\noalign{}
\endlastfoot
electief & 69109 & 2.2 & -4.9 & 43 & 19.5 \\
niet-electief & 286 & -2.6 & -7.3 & 31.4 & -12.1 \\
\end{longtable}

The mean relative deviation between planned and observed durations
remains broadly stable over the observed years. In 2022 the average
deviation is 0.1 percent, increasing slightly to 1.9 percent in 2023. A
clearer rise appears in 2024, reaching 4.8 percent, while the partial
data for 2025 show a decrease to 2.1 percent.

This pattern indicates that planning accuracy at Lanaken has remained
relatively consistent over time, with only moderate year to year
fluctuations. The differences are small and likely reflect normal
operational variation rather than systematic changes in planning
quality. Overall, the stability of the yearly averages suggests that the
underlying scheduling process is well established, even though
individual level variability continues to affect day to day accuracy.
Figure \ref{fig:reldev_year} illustrates these trends.

\begin{figure}[H]

{\centering \includegraphics[width=0.7\linewidth,]{In-Depth_Analysis_Lanaken_files/figure-latex/rel_duration_dev_year-1} 

}

\caption{Relative duration deviation by year \label{fig:reldev_year}}\label{fig:rel_duration_dev_year}
\end{figure}

The relative deviation between planned and observed durations at Campus
Lanaken shows a consistent pattern across the regular workweek. Mean
deviations range from 0.8 percent to 3.9 percent, indicating that
planning accuracy is generally stable from Monday through Friday. The
slightly higher deviation on Tuesday, followed by a moderate value on
Wednesday, likely reflects variation in the types of procedures
performed rather than structural scheduling differences. Lower values on
Thursday and Friday suggest that case mix later in the week may be more
predictable. Overall, the figure shows that at Lanaken the alignment
between planned and observed durations is steady across weekdays, with
only modest fluctuations tied to day specific activity profiles.

\begin{figure}[H]

{\centering \includegraphics[width=0.7\linewidth,]{In-Depth_Analysis_Lanaken_files/figure-latex/rel_dev_weekday-1} 

}

\caption{Relative deviation by weekday \label{fig:reldev_weekday}}\label{fig:rel_dev_weekday}
\end{figure}

Variation in deviation differs considerably across operating rooms. The
largest average deviations are concentrated in the complex surgery rooms
of the LO block, most notably LO03, where the mean absolute deviation is
16.4 minutes and 45.9 percent of cases run longer than planned. Other
high variability rooms such as LO05, LO01, and LO04 also show average
deviations between 10.6 and 13.5 minutes, reflecting the higher
uncertainty typical of lengthy or technically demanding procedures.

By contrast, rooms such as LP01 and LP02, mainly associated with short
and standardized interventions, display mean deviations below 6 minutes
and a roughly balanced share of shorter and longer cases. These
differences illustrate how the degree of procedural complexity and
specialization strongly influences planning precision. Table
\ref{tab:ordeviation_or} provides the detailed overview.

\begin{longtable}[]{@{}
  >{\centering\arraybackslash}p{(\linewidth - 12\tabcolsep) * \real{0.0917}}
  >{\centering\arraybackslash}p{(\linewidth - 12\tabcolsep) * \real{0.0667}}
  >{\centering\arraybackslash}p{(\linewidth - 12\tabcolsep) * \real{0.1250}}
  >{\centering\arraybackslash}p{(\linewidth - 12\tabcolsep) * \real{0.1083}}
  >{\centering\arraybackslash}p{(\linewidth - 12\tabcolsep) * \real{0.2417}}
  >{\centering\arraybackslash}p{(\linewidth - 12\tabcolsep) * \real{0.1167}}
  >{\centering\arraybackslash}p{(\linewidth - 12\tabcolsep) * \real{0.2500}}@{}}
\caption{Duration deviations per operating room, ordered by mean
absolute deviation \label{tab:ordeviation_or}}\tabularnewline
\toprule\noalign{}
\begin{minipage}[b]{\linewidth}\centering
ActualOR
\end{minipage} & \begin{minipage}[b]{\linewidth}\centering
n
\end{minipage} & \begin{minipage}[b]{\linewidth}\centering
Mean abs dev
\end{minipage} & \begin{minipage}[b]{\linewidth}\centering
Longer (\%)
\end{minipage} & \begin{minipage}[b]{\linewidth}\centering
Mean relative dev (longer)
\end{minipage} & \begin{minipage}[b]{\linewidth}\centering
Shorter (\%)
\end{minipage} & \begin{minipage}[b]{\linewidth}\centering
Mean relative dev (shorter)
\end{minipage} \\
\midrule\noalign{}
\endfirsthead
\toprule\noalign{}
\begin{minipage}[b]{\linewidth}\centering
ActualOR
\end{minipage} & \begin{minipage}[b]{\linewidth}\centering
n
\end{minipage} & \begin{minipage}[b]{\linewidth}\centering
Mean abs dev
\end{minipage} & \begin{minipage}[b]{\linewidth}\centering
Longer (\%)
\end{minipage} & \begin{minipage}[b]{\linewidth}\centering
Mean relative dev (longer)
\end{minipage} & \begin{minipage}[b]{\linewidth}\centering
Shorter (\%)
\end{minipage} & \begin{minipage}[b]{\linewidth}\centering
Mean relative dev (shorter)
\end{minipage} \\
\midrule\noalign{}
\endhead
\bottomrule\noalign{}
\endlastfoot
LO03 & 4196 & 16.4 & 45.9 & 29.3 & 51.1 & -21.2 \\
LO05 & 5206 & 13.5 & 41.2 & 30.8 & 55.2 & -23.4 \\
LO01 & 5202 & 12.3 & 39.2 & 27.6 & 57.5 & -21 \\
LO04 & 6009 & 10.6 & 45.1 & 30.7 & 50.8 & -20.9 \\
LO02 & 9974 & 7.8 & 41.5 & 30.1 & 52.9 & -20.7 \\
LP02 & 18619 & 5.6 & 38.6 & 43.7 & 54.8 & -31.2 \\
LP01 & 20189 & 4.8 & 43.6 & 46.3 & 48.8 & -29.5 \\
\end{longtable}

Among the most extreme deviations in surgery duration, longer and
shorter cases occur in almost equal proportions. Around 66.8 percent of
the top one percent of deviations correspond to procedures that last
longer than planned, while 33.2 percent finish earlier than expected.
This contrasts with earlier observations based on a different dataset,
where overruns dominated more clearly.

Although both directions are represented among the extreme cases, their
operational impact remains different. Early finishes generally create
limited idle time, whereas overruns can delay subsequent cases and
increase the likelihood of overtime. The near balance between longer and
shorter cases in the current extreme deviations suggests that large
underestimations and overestimations both contribute meaningfully to
schedule variability, reinforcing the need to monitor procedures prone
to substantial deviations in either direction. Table
\ref{tab:extreme_dir} summarizes these results.

\begin{longtable}[]{@{}
  >{\centering\arraybackslash}p{(\linewidth - 4\tabcolsep) * \real{0.3194}}
  >{\centering\arraybackslash}p{(\linewidth - 4\tabcolsep) * \real{0.0833}}
  >{\centering\arraybackslash}p{(\linewidth - 4\tabcolsep) * \real{0.1806}}@{}}
\caption{Direction of deviation among the 1\% most extreme cases
(Lanaken) \label{tab:extreme_dir}}\tabularnewline
\toprule\noalign{}
\begin{minipage}[b]{\linewidth}\centering
Direction
\end{minipage} & \begin{minipage}[b]{\linewidth}\centering
n
\end{minipage} & \begin{minipage}[b]{\linewidth}\centering
Percentage
\end{minipage} \\
\midrule\noalign{}
\endfirsthead
\toprule\noalign{}
\begin{minipage}[b]{\linewidth}\centering
Direction
\end{minipage} & \begin{minipage}[b]{\linewidth}\centering
n
\end{minipage} & \begin{minipage}[b]{\linewidth}\centering
Percentage
\end{minipage} \\
\midrule\noalign{}
\endhead
\bottomrule\noalign{}
\endlastfoot
Longer than planned & 474 & 66.8\% \\
Shorter than planned & 236 & 33.2\% \\
\end{longtable}

\subsection{A closer look on the coefficient of
variation}\label{a-closer-look-on-the-coefficient-of-variation}

The monthly coefficient of variation for planning deviation displays a
highly volatile pattern, as shown in Figure \ref{fig:cv_planning}. Most
months cluster around relatively low values between 1.0 and 1.5, but
these stable periods are repeatedly interrupted by sharp, isolated
spikes. Early 2022 shows modest variation until a sudden jump to nearly
4, followed shortly by another month above 3. Similar single-month
surges reappear throughout 2023 and 2024, with values that briefly
exceed 3 before dropping back to the lower baseline.

These spikes are sporadic rather than sustained. They do not form a
rising or falling trend but instead appear as short bursts of unusually
high proportional deviation, likely driven by temporary shifts in case
mix or isolated outliers. By contrast, 2025 exhibits much steadier
behaviour, with the coefficient remaining close to 1.2 and no visible
extreme peaks.

Overall, the pattern is characterised by a low, stable baseline
punctuated by occasional large fluctuations. This indicates that
planning accuracy at Lanaken is generally consistent, but certain months
experience markedly higher relative dispersion, reflecting short-term
variability rather than a structural change over time.

\begin{figure}[H]

{\centering \includegraphics[width=0.7\linewidth,]{In-Depth_Analysis_Lanaken_files/figure-latex/cv_monthly-1} 

}

\caption{Monthly evolution of coefficient of variation (Planning deviation) \label{fig:cv_planning}}\label{fig:cv_monthly}
\end{figure}

\subsubsection{Variability by planned duration
group}\label{variability-by-planned-duration-group}

Variation in surgical timing can be examined from two complementary
perspectives that describe different aspects of performance. The first
is the coefficient of variation of observed duration, calculated as the
standard deviation divided by the mean of the observed surgical times.
This indicator expresses how dispersed the recorded durations are
relative to their average. In a dataset such as this one, which contains
a large number of both short and long operations, a high coefficient of
variation mainly reflects the broad spread in case duration across
procedure types rather than instability within individual surgeries.

The second perspective is the coefficient of variation of the planning
deviation, which measures the consistency of planning accuracy. It is
computed as the standard deviation of the absolute deviation between
planned and observed time, divided by the mean of that deviation. This
metric focuses on how variable the alignment between planned and
realised durations is, regardless of how long the surgeries themselves
last.

Comparing both perspectives helps disentangle two components of
variability: the structural heterogeneity in procedure duration and the
variability in how precisely these durations are planned. Understanding
how these two aspects evolve provides insight into whether deviations
primarily originate from the diversity of surgical activity or from
variation in scheduling accuracy. The following table summarises these
indicators across five duration categories and serves as a baseline for
the more detailed, category-specific analyses presented in the
subsequent subsections. All values are reported in Table
\ref{tab:cv_group}, which provides the exact CV measures for both
observed durations and planning deviations.

The summary statistics confirm a clear scaling effect in relative
variability. The coefficient of variation of observed durations
decreases as planned duration increases: procedures planned for less
than 30 minutes have a CV of 0.5210286, procedures in the 31--60 minute
group show 0.3819727, the 61--90 minute group 0.3449165, the 91--180
minute group 0.3180666, and the longest category over 180 minutes shows
0.6921203 This pattern is consistent with the fact that shorter cases
amplify small absolute fluctuations when expressed relative to their
mean, while longer cases distribute absolute variation over a broader
time base.

For the planning deviation, coefficients of variation are considerably
higher across all groups, reflecting the larger proportional spread in
the differences between planned and actual duration. These values range
from 0.9900561 in the shortest category to 0.9127430 in the 31--60
minute group, 0.8386945 in the 61--90 minute group, 0.8976667 in the
91--180 minute group, and 0.9899863 in procedures lasting more than 180
minutes. This pattern suggests that longer and more complex procedures,
while not necessarily less stable in execution, exhibit greater
variation in planning accuracy.

Together, these figures show that relative variability in observed
duration is largely structural and reflects case mix composition,
whereas variability in planning deviation captures the irregularity in
scheduling precision. The distinction between the two metrics provides
the foundation for the next analyses, which explore how these patterns
manifest within each duration category and across operating rooms.

\begin{table}

\caption{\label{tab:cv group}Summary of mean duration, SD, and CV by planned duration group (Observed vs. Planning deviation) \label{tab:cv_group}}
\centering
\begin{tabular}[t]{lrrrrrrr}
\toprule
Group & n & Mean duration & SD duration & CV duration & Mean planning & SD planning & CV planning\\
\midrule
<30 & 50060 & 17.78666 & 9.267356 & 0.5210286 & 5.337355 & 5.284281 & 0.9900561\\
31–60 & 12315 & 43.15436 & 16.483787 & 0.3819727 & 10.223386 & 9.331324 & 0.9127430\\
61–90 & 4320 & 75.02523 & 25.877444 & 0.3449165 & 18.060185 & 15.146979 & 0.8386945\\
91–180 & 2662 & 108.92186 & 34.644410 & 0.3180666 & 23.638618 & 21.219601 & 0.8976667\\
> 180 & 38 & 109.57895 & 75.841810 & 0.6921203 & 307.578947 & 304.498936 & 0.9899863\\
\bottomrule
\end{tabular}
\end{table}

The monthly evolution of the coefficient of variation for observed
durations shows a highly stable pattern for four of the five planned
duration groups, as visible in Figure \ref{fig:cv_obs_group}. The
\textless30, 31--60, 61--90, and 91--180 minute groups all form dense,
nearly flat trajectories over time, with only minor month-to-month
fluctuations. These curves remain tightly clustered: the \textless30
minute group sits consistently around 0.4 to 0.5, whereas the three
mid-length groups remain lower, typically between 0.25 and 0.4. The
clear separation between short procedures and the mid-range groups
persists throughout the entire period, reflecting stable proportional
variability.

In contrast, the \textgreater180 minute group behaves very differently.
Instead of a continuous line, it appears only sporadically and shows
extreme swings, from near zero to values close to 1.0. These abrupt
jumps reflect the small number of very long procedures performed each
month, where even a few cases can cause large shifts in the coefficient
of variation. The volatility therefore stems from data sparsity rather
than a structural change in duration variability.

Overall, the figure reveals a stable and well-behaved pattern for the
majority of procedure lengths, while the longest-duration category
displays irregular and highly sensitive behaviour due to low monthly
counts.

\begin{figure}[H]

{\centering \includegraphics[width=0.7\linewidth,]{In-Depth_Analysis_Lanaken_files/figure-latex/cv obs group-1} 

}

\caption{Monthly evolution of coefficient of variation (Observed durations) \label{fig:cv_obs_group}}\label{fig:cv obs group}
\end{figure}

The coefficient of variation of planning deviation shows a dynamic
pattern over time, with marked differences between duration groups, as
illustrated in Figure \ref{fig:cv_planning_group}. Across most months,
the mid-range categories (31--60, 61--90, and 91--180 minutes) form a
compact cluster, fluctuating between roughly 0.7 and 1.1. Their tight
trajectories indicate that planning deviations for these procedures are
comparatively stable and change only modestly from month to month.

Short procedures (\textless30 minutes) display noticeably larger swings.
Their curve oscillates more widely than the mid-range groups and reaches
several pronounced peaks around 1.2 to 1.3. These spikes suggest periods
in which small absolute deviations become proportionally large for very
short cases, making the coefficient of variation more sensitive to case
mix shifts.

The \textgreater180 minute group behaves quite differently: instead of a
continuous trajectory, it appears intermittently with large jumps,
sometimes falling below 0.6 and at other times rising above 1.3. These
swings are driven by small monthly sample sizes, where just a few
unusually short or long deviations can cause sharp changes in the
coefficient.

Overall, the figure shows that planning deviation is most stable in the
mid-duration range, more volatile for short cases, and highly irregular
for the longest procedures because of sparse data.

\begin{figure}[H]

{\centering \includegraphics[width=0.7\linewidth,]{In-Depth_Analysis_Lanaken_files/figure-latex/cv planning group-1} 

}

\caption{Monthly evolution of coefficient of variation (Planning deviation) \label{fig:cv_planning_group}}\label{fig:cv planning group}
\end{figure}

Together, these patterns indicate that the surgical process itself is
proportionally consistent, while planning accuracy is more sensitive to
case mix shifts, episodic complexity, and the inherent difficulty of
predicting very short and very long procedures. The next sections
examine these dynamics in more detail by analysing each duration group
separately and comparing the distribution of coefficients of variation
across operating rooms.

\vspace{1em}

\noindent\textbf{Surgeries planned under 30 minutes} \vspace{0.5em}

Short procedures, typically representing high volume and standardized
interventions, show moderate relative variability across operating
rooms. For this category, coefficients of variation of observed
durations range roughly between 0.35 and 0.63, which reflects the
proportional impact of small absolute timing fluctuations on short
cases. These values are expected for brief procedures, where even
limited variation translates into relatively large coefficients once
expressed relative to their mean duration. Table \ref{tab:cv_room_lt30}
summarises these indicators for all rooms with at least eleven surgeries
in this duration group.

The distribution across rooms suggests a broadly homogeneous pattern.
Units such as LP01, LP02, and LO05 show higher CV values between 0.45
and 0.62, while others including LO02, LO04, LO01, and LO03 show lower
variability, with values between 0.35 and 0.36. These differences are
small in absolute terms and mainly reflect differences in the mix of
short procedures performed in each room rather than substantive
disparities in timing consistency.

For the coefficient of variation of planning deviation, values are
systematically higher and span a wider range, with most rooms falling
between approximately 0.9127 and 1.213 This confirms that relative
fluctuations in planning accuracy exceed those associated with the
observed durations themselves. The lowest CV deviation values appear in
rooms such as LP02, LO01, and LO04, while higher values are found in
rooms including LO02, LO05, and LO03

Overall, short procedures are characterized by consistent execution
times and broadly comparable planning accuracy across operating rooms.
The relatively high coefficients primarily reflect proportional scaling
effects inherent to short interventions rather than substantial
operational variability. Subsequent sections examine whether similar
patterns hold for procedures of longer planned duration, where absolute
variability becomes more dominant and planning precision potentially
more challenging.

\begin{longtable}[]{@{}
  >{\centering\arraybackslash}p{(\linewidth - 12\tabcolsep) * \real{0.0864}}
  >{\centering\arraybackslash}p{(\linewidth - 12\tabcolsep) * \real{0.0988}}
  >{\centering\arraybackslash}p{(\linewidth - 12\tabcolsep) * \real{0.1975}}
  >{\centering\arraybackslash}p{(\linewidth - 12\tabcolsep) * \real{0.1728}}
  >{\centering\arraybackslash}p{(\linewidth - 12\tabcolsep) * \real{0.1728}}
  >{\centering\arraybackslash}p{(\linewidth - 12\tabcolsep) * \real{0.1358}}
  >{\centering\arraybackslash}p{(\linewidth - 12\tabcolsep) * \real{0.1358}}@{}}
\caption{Coefficient of variation per operating room for surgeries
planned \textless30 minutes (only rooms with \textgreater10 surgeries
shown) \label{tab:cv_room_lt30} (continued below)}\tabularnewline
\toprule\noalign{}
\begin{minipage}[b]{\linewidth}\centering
OR
\end{minipage} & \begin{minipage}[b]{\linewidth}\centering
n
\end{minipage} & \begin{minipage}[b]{\linewidth}\centering
Mean duration
\end{minipage} & \begin{minipage}[b]{\linewidth}\centering
SD duration
\end{minipage} & \begin{minipage}[b]{\linewidth}\centering
CV duration
\end{minipage} & \begin{minipage}[b]{\linewidth}\centering
Mean dev
\end{minipage} & \begin{minipage}[b]{\linewidth}\centering
SD dev
\end{minipage} \\
\midrule\noalign{}
\endfirsthead
\toprule\noalign{}
\begin{minipage}[b]{\linewidth}\centering
OR
\end{minipage} & \begin{minipage}[b]{\linewidth}\centering
n
\end{minipage} & \begin{minipage}[b]{\linewidth}\centering
Mean duration
\end{minipage} & \begin{minipage}[b]{\linewidth}\centering
SD duration
\end{minipage} & \begin{minipage}[b]{\linewidth}\centering
CV duration
\end{minipage} & \begin{minipage}[b]{\linewidth}\centering
Mean dev
\end{minipage} & \begin{minipage}[b]{\linewidth}\centering
SD dev
\end{minipage} \\
\midrule\noalign{}
\endhead
\bottomrule\noalign{}
\endlastfoot
LP01 & 20114 & 14.73 & 6.893 & 0.468 & 4.787 & 4.664 \\
LP02 & 18090 & 15.61 & 7.157 & 0.4585 & 5.157 & 4.707 \\
LO02 & 6504 & 25.09 & 8.913 & 0.3552 & 5.905 & 5.948 \\
LO04 & 2714 & 29.12 & 10.28 & 0.3529 & 7.336 & 7.069 \\
LO05 & 1291 & 22.26 & 13.96 & 0.6269 & 6.697 & 8.123 \\
LO03 & 707 & 29.84 & 10.52 & 0.3525 & 7.395 & 7.17 \\
LO01 & 640 & 30.76 & 11.04 & 0.3589 & 8.483 & 8.032 \\
\end{longtable}

\begin{longtable}[]{@{}
  >{\centering\arraybackslash}p{(\linewidth - 0\tabcolsep) * \real{0.1250}}@{}}
\toprule\noalign{}
\begin{minipage}[b]{\linewidth}\centering
CV dev
\end{minipage} \\
\midrule\noalign{}
\endhead
\bottomrule\noalign{}
\endlastfoot
0.9743 \\
0.9127 \\
1.007 \\
0.9636 \\
1.213 \\
0.9696 \\
0.9469 \\
\end{longtable}

\vspace{1em}

\noindent\textbf{Surgeries planned 31–60 minutes} \vspace{0.5em}

Procedures with a planned duration of less than 30 minutes show
relatively high proportional variability across operating rooms,
reflecting the sensitivity of short procedures to even small absolute
fluctuations. For this category, coefficients of variation of observed
durations range from about 0.35 to 0.63, as shown in Table
\ref{tab:cv_room_under30}. The lowest values appear in LO02 (0.3552),
LO04 (0.3529), and LO01 (0.3589), while rooms such as LP01 (0.468) and
LP02 (0.4585) fall in the mid-range. LO05 shows the highest CV\_duration
(0.6269), which likely reflects a more heterogeneous short-procedure mix
rather than operational inconsistencies.

For the coefficient of variation of planning deviation, values are again
higher across all rooms, ranging from about 0.91 to 1.21. As shown in
Table \ref{tab:cv_room_under30}, the lowest CV\_deviation values appear
in LP02 (0.9127), LO01 (0.9469), and LO04 (0.9636), while LO03 (0.9669)
and LP01 (0.9743) lie near the centre of the distribution. The highest
values occur in LO02 (1.007) and LO05 (1.213), indicating that planning
deviations for the shortest procedures often match or exceed their mean,
which is expected because even a few minutes of difference represent a
large proportional deviation.

Taken together, these results show that procedures planned for less than
30 minutes display predictable execution in absolute terms, but
inherently high relative variability in both duration and planning
deviation. The observed coefficients largely reflect the scaling effects
typical of very short procedures rather than meaningful operational
differences between rooms.

\begin{longtable}[]{@{}
  >{\centering\arraybackslash}p{(\linewidth - 14\tabcolsep) * \real{0.0805}}
  >{\centering\arraybackslash}p{(\linewidth - 14\tabcolsep) * \real{0.0805}}
  >{\centering\arraybackslash}p{(\linewidth - 14\tabcolsep) * \real{0.1839}}
  >{\centering\arraybackslash}p{(\linewidth - 14\tabcolsep) * \real{0.1609}}
  >{\centering\arraybackslash}p{(\linewidth - 14\tabcolsep) * \real{0.1609}}
  >{\centering\arraybackslash}p{(\linewidth - 14\tabcolsep) * \real{0.1264}}
  >{\centering\arraybackslash}p{(\linewidth - 14\tabcolsep) * \real{0.1034}}
  >{\centering\arraybackslash}p{(\linewidth - 14\tabcolsep) * \real{0.1034}}@{}}
\caption{Coefficient of variation per operating room for surgeries
planned 31--60 minutes (only rooms with \textgreater10 surgeries shown)
\label{tab:cv_room_31_60}}\tabularnewline
\toprule\noalign{}
\begin{minipage}[b]{\linewidth}\centering
OR
\end{minipage} & \begin{minipage}[b]{\linewidth}\centering
n
\end{minipage} & \begin{minipage}[b]{\linewidth}\centering
Mean duration
\end{minipage} & \begin{minipage}[b]{\linewidth}\centering
SD duration
\end{minipage} & \begin{minipage}[b]{\linewidth}\centering
CV duration
\end{minipage} & \begin{minipage}[b]{\linewidth}\centering
Mean dev
\end{minipage} & \begin{minipage}[b]{\linewidth}\centering
SD dev
\end{minipage} & \begin{minipage}[b]{\linewidth}\centering
CV dev
\end{minipage} \\
\midrule\noalign{}
\endfirsthead
\toprule\noalign{}
\begin{minipage}[b]{\linewidth}\centering
OR
\end{minipage} & \begin{minipage}[b]{\linewidth}\centering
n
\end{minipage} & \begin{minipage}[b]{\linewidth}\centering
Mean duration
\end{minipage} & \begin{minipage}[b]{\linewidth}\centering
SD duration
\end{minipage} & \begin{minipage}[b]{\linewidth}\centering
CV duration
\end{minipage} & \begin{minipage}[b]{\linewidth}\centering
Mean dev
\end{minipage} & \begin{minipage}[b]{\linewidth}\centering
SD dev
\end{minipage} & \begin{minipage}[b]{\linewidth}\centering
CV dev
\end{minipage} \\
\midrule\noalign{}
\endhead
\bottomrule\noalign{}
\endlastfoot
LO02 & 3085 & 46.58 & 16.21 & 0.348 & 10.08 & 8.62 & 0.8553 \\
LO01 & 2915 & 44.35 & 15.1 & 0.3406 & 9.718 & 8.979 & 0.9239 \\
LO05 & 2175 & 41.6 & 15.68 & 0.377 & 9.789 & 9.248 & 0.9448 \\
LO04 & 2172 & 41.38 & 16.62 & 0.4017 & 9.151 & 8.786 & 0.9602 \\
LO03 & 1467 & 43.8 & 17.72 & 0.4045 & 11.2 & 10.98 & 0.9803 \\
LP02 & 437 & 27.64 & 14.72 & 0.5326 & 17.72 & 9.466 & 0.5341 \\
LP01 & 64 & 28.38 & 12.31 & 0.434 & 17.88 & 9.297 & 0.5201 \\
\end{longtable}

\vspace{1em}

\noindent\textbf{Surgeries planned 61–90 minutes} \vspace{0.5em}

Procedures with a planned duration of 31--60 minutes show moderate
relative variability across operating rooms. For this category,
coefficients of variation of observed durations cluster between about
0.34 and 0.53, as shown in Table \ref{tab:cv_room_31_60}. The lowest
values appear in LO01 (0.3406) and LO02 (0.348), followed closely by
LO04 (0.4017) and LO03 (0.4045). Higher variability is observed in LO05
(0.377) and LP01 (0.434), while LP02 shows the highest CV\_duration at
0.5326. These differences likely reflect variation in the underlying
case mix rather than structural workflow differences between rooms.

For the coefficient of variation of planning deviation, values are again
higher across all rooms, ranging from approximately 0.52 to 0.98. The
lowest CV\_dev values appear in LP01 (0.5201) and LP02 (0.6341), while
rooms such as LO01 (0.9239), LO05 (0.9448), LO03 (0.9803), and LO04
(0.9602) display substantially higher dispersion. LO02 (0.8553) lies in
the mid-range. These patterns indicate that relative variation in
planning accuracy generally exceeds the variation in the durations
themselves, especially in the general-purpose rooms where procedure
heterogeneity is greater.

Taken together, these results show that procedures planned for 31--60
minutes display consistent execution times and moderate variation in
planning accuracy across operating rooms. The observed coefficients
largely follow proportional scaling effects typical of medium-length
procedures rather than substantial operational differences.

\begin{longtable}[]{@{}
  >{\centering\arraybackslash}p{(\linewidth - 14\tabcolsep) * \real{0.0805}}
  >{\centering\arraybackslash}p{(\linewidth - 14\tabcolsep) * \real{0.0805}}
  >{\centering\arraybackslash}p{(\linewidth - 14\tabcolsep) * \real{0.1839}}
  >{\centering\arraybackslash}p{(\linewidth - 14\tabcolsep) * \real{0.1609}}
  >{\centering\arraybackslash}p{(\linewidth - 14\tabcolsep) * \real{0.1609}}
  >{\centering\arraybackslash}p{(\linewidth - 14\tabcolsep) * \real{0.1264}}
  >{\centering\arraybackslash}p{(\linewidth - 14\tabcolsep) * \real{0.1034}}
  >{\centering\arraybackslash}p{(\linewidth - 14\tabcolsep) * \real{0.1034}}@{}}
\caption{Coefficient of variation per operating room for surgeries
planned 61--90 minutes (only rooms with \textgreater10 surgeries shown)
\label{tab:cv_room_61_90}}\tabularnewline
\toprule\noalign{}
\begin{minipage}[b]{\linewidth}\centering
OR
\end{minipage} & \begin{minipage}[b]{\linewidth}\centering
n
\end{minipage} & \begin{minipage}[b]{\linewidth}\centering
Mean duration
\end{minipage} & \begin{minipage}[b]{\linewidth}\centering
SD duration
\end{minipage} & \begin{minipage}[b]{\linewidth}\centering
CV duration
\end{minipage} & \begin{minipage}[b]{\linewidth}\centering
Mean dev
\end{minipage} & \begin{minipage}[b]{\linewidth}\centering
SD dev
\end{minipage} & \begin{minipage}[b]{\linewidth}\centering
CV dev
\end{minipage} \\
\midrule\noalign{}
\endfirsthead
\toprule\noalign{}
\begin{minipage}[b]{\linewidth}\centering
OR
\end{minipage} & \begin{minipage}[b]{\linewidth}\centering
n
\end{minipage} & \begin{minipage}[b]{\linewidth}\centering
Mean duration
\end{minipage} & \begin{minipage}[b]{\linewidth}\centering
SD duration
\end{minipage} & \begin{minipage}[b]{\linewidth}\centering
CV duration
\end{minipage} & \begin{minipage}[b]{\linewidth}\centering
Mean dev
\end{minipage} & \begin{minipage}[b]{\linewidth}\centering
SD dev
\end{minipage} & \begin{minipage}[b]{\linewidth}\centering
CV dev
\end{minipage} \\
\midrule\noalign{}
\endhead
\bottomrule\noalign{}
\endlastfoot
LO01 & 1312 & 67.88 & 20.3 & 0.2991 & 15.44 & 12.2 & 0.7901 \\
LO05 & 1154 & 79.28 & 24.02 & 0.3029 & 18.96 & 14.74 & 0.7777 \\
LO03 & 759 & 84.53 & 27.85 & 0.3295 & 19.81 & 16.72 & 0.8441 \\
LO04 & 639 & 83.87 & 29.41 & 0.3506 & 19.53 & 18.92 & 0.969 \\
LO02 & 354 & 59.84 & 20.3 & 0.3392 & 15.75 & 13.66 & 0.8672 \\
LP02 & 92 & 45.33 & 15.87 & 0.3502 & 27.8 & 12.36 & 0.4444 \\
\end{longtable}

\vspace{1em}

\noindent\textbf{Surgeries planned 91–180 minutes} \vspace{0.5em}

Surgeries with a planned duration between 91 and 180 minutes include a
broad mix of major inpatient operations and complex elective cases.
Within this category, the coefficient of variation of observed durations
is moderate, ranging mostly between 0.27 and 0.41 across operating
rooms, as shown in Table \ref{tab:cv_room_91_180}. The lowest CV values
are observed in LO04 (0.2768) and LO03 (0.2909), indicating a tightly
clustered distribution of operative times. Slightly higher values appear
in LO05 (0.3322) and LO02 (0.3759), while LO01 shows the highest
proportional variability in this group with a CV of 0.4078. These
differences remain limited and suggest proportionally consistent
execution times despite the broader absolute duration range associated
with longer procedures.

The coefficient of variation of planning deviation shows a similarly
narrow dispersion. Most rooms fall between 0.76 and 1.01, with the
lowest values seen in LO02 (0.7677) and LO04 (0.8273). LO05 (0.8459) and
LO03 (0.8973) occupy the mid-range, while LO01 shows the largest
relative dispersion in planning deviation (1.005). Despite these
differences, the overall spread remains constrained, indicating stable
planning accuracy across rooms even as case mix becomes more
heterogeneous.

Across operating rooms, both indicators show restrained variability,
implying that once surgeries reach this duration range, timing
performance becomes comparatively stable. In contrast to shorter
procedures, proportional variation in duration and planning deviation
becomes smaller relative to total operative time. These mid- to
long-duration operations therefore represent one of the most consistent
segments of the surgical program in terms of timing reliability.

\begin{longtable}[]{@{}
  >{\centering\arraybackslash}p{(\linewidth - 14\tabcolsep) * \real{0.0805}}
  >{\centering\arraybackslash}p{(\linewidth - 14\tabcolsep) * \real{0.0805}}
  >{\centering\arraybackslash}p{(\linewidth - 14\tabcolsep) * \real{0.1839}}
  >{\centering\arraybackslash}p{(\linewidth - 14\tabcolsep) * \real{0.1609}}
  >{\centering\arraybackslash}p{(\linewidth - 14\tabcolsep) * \real{0.1609}}
  >{\centering\arraybackslash}p{(\linewidth - 14\tabcolsep) * \real{0.1264}}
  >{\centering\arraybackslash}p{(\linewidth - 14\tabcolsep) * \real{0.1034}}
  >{\centering\arraybackslash}p{(\linewidth - 14\tabcolsep) * \real{0.1034}}@{}}
\caption{Coefficient of variation per operating room for surgeries
planned 91--180 minutes (only rooms with \textgreater10 surgeries shown)
\label{tab:cv_room_91_180}}\tabularnewline
\toprule\noalign{}
\begin{minipage}[b]{\linewidth}\centering
OR
\end{minipage} & \begin{minipage}[b]{\linewidth}\centering
n
\end{minipage} & \begin{minipage}[b]{\linewidth}\centering
Mean duration
\end{minipage} & \begin{minipage}[b]{\linewidth}\centering
SD duration
\end{minipage} & \begin{minipage}[b]{\linewidth}\centering
CV duration
\end{minipage} & \begin{minipage}[b]{\linewidth}\centering
Mean dev
\end{minipage} & \begin{minipage}[b]{\linewidth}\centering
SD dev
\end{minipage} & \begin{minipage}[b]{\linewidth}\centering
CV dev
\end{minipage} \\
\midrule\noalign{}
\endfirsthead
\toprule\noalign{}
\begin{minipage}[b]{\linewidth}\centering
OR
\end{minipage} & \begin{minipage}[b]{\linewidth}\centering
n
\end{minipage} & \begin{minipage}[b]{\linewidth}\centering
Mean duration
\end{minipage} & \begin{minipage}[b]{\linewidth}\centering
SD duration
\end{minipage} & \begin{minipage}[b]{\linewidth}\centering
CV duration
\end{minipage} & \begin{minipage}[b]{\linewidth}\centering
Mean dev
\end{minipage} & \begin{minipage}[b]{\linewidth}\centering
SD dev
\end{minipage} & \begin{minipage}[b]{\linewidth}\centering
CV dev
\end{minipage} \\
\midrule\noalign{}
\endhead
\bottomrule\noalign{}
\endlastfoot
LO03 & 1251 & 108.9 & 31.68 & 0.2909 & 22.1 & 19.83 & 0.8973 \\
LO05 & 573 & 109.1 & 36.23 & 0.3322 & 25.3 & 21.4 & 0.8459 \\
LO04 & 482 & 105.4 & 29.17 & 0.2768 & 22.67 & 18.75 & 0.8273 \\
LO01 & 326 & 114.5 & 46.68 & 0.4078 & 26.8 & 26.94 & 1.005 \\
LO02 & 29 & 105.8 & 39.78 & 0.3759 & 34.79 & 26.71 & 0.7677 \\
\end{longtable}

\vspace{1em}

\noindent\textbf{Surgeries planned over 180 minutes} \vspace{0.5em}

Procedures with a planned duration exceeding three hours represent the
most complex and resource intensive part of the surgical program. Within
this group, the coefficient of variation of observed durations remains
moderate, ranging from 0.67 to 0.80 across operating rooms, as shown in
Table \ref{tab:cv_room_over180}. The lowest CV value appears in LO03
(0.6698), while LO05 shows the highest proportional variability at
0.7981. Interpretation remains cautious because both rooms have
relatively small case counts, which can amplify variability in these
proportional measures.

The coefficient of variation of planning deviation is generally close to
or above one. The rooms fall between 0.8391 and 1.027, with LO03 showing
the lowest value (0.8391) and LO05 the highest (1.027). These
differences reflect the wide absolute spread in deviations that
naturally accompanies long and complex procedures rather than systematic
inconsistencies in planning accuracy.

Overall, surgeries longer than three hours display stable and
proportionally uniform duration patterns across the operating complex.
Observed durations show limited relative variability, and planning
deviations exhibit consistent proportional spread. The higher
coefficients observed in a few rooms likely result from the
concentration of rare or high complexity procedures rather than
persistent scheduling issues.

\begin{longtable}[]{@{}
  >{\centering\arraybackslash}p{(\linewidth - 14\tabcolsep) * \real{0.0824}}
  >{\centering\arraybackslash}p{(\linewidth - 14\tabcolsep) * \real{0.0588}}
  >{\centering\arraybackslash}p{(\linewidth - 14\tabcolsep) * \real{0.1882}}
  >{\centering\arraybackslash}p{(\linewidth - 14\tabcolsep) * \real{0.1647}}
  >{\centering\arraybackslash}p{(\linewidth - 14\tabcolsep) * \real{0.1647}}
  >{\centering\arraybackslash}p{(\linewidth - 14\tabcolsep) * \real{0.1294}}
  >{\centering\arraybackslash}p{(\linewidth - 14\tabcolsep) * \real{0.1059}}
  >{\centering\arraybackslash}p{(\linewidth - 14\tabcolsep) * \real{0.1059}}@{}}
\caption{Coefficient of variation per operating room for surgeries
planned \textgreater{} 180 minutes (only rooms with \textgreater10
surgeries shown) \label{tab:cv_room_over180}}\tabularnewline
\toprule\noalign{}
\begin{minipage}[b]{\linewidth}\centering
OR
\end{minipage} & \begin{minipage}[b]{\linewidth}\centering
n
\end{minipage} & \begin{minipage}[b]{\linewidth}\centering
Mean duration
\end{minipage} & \begin{minipage}[b]{\linewidth}\centering
SD duration
\end{minipage} & \begin{minipage}[b]{\linewidth}\centering
CV duration
\end{minipage} & \begin{minipage}[b]{\linewidth}\centering
Mean dev
\end{minipage} & \begin{minipage}[b]{\linewidth}\centering
SD dev
\end{minipage} & \begin{minipage}[b]{\linewidth}\centering
CV dev
\end{minipage} \\
\midrule\noalign{}
\endfirsthead
\toprule\noalign{}
\begin{minipage}[b]{\linewidth}\centering
OR
\end{minipage} & \begin{minipage}[b]{\linewidth}\centering
n
\end{minipage} & \begin{minipage}[b]{\linewidth}\centering
Mean duration
\end{minipage} & \begin{minipage}[b]{\linewidth}\centering
SD duration
\end{minipage} & \begin{minipage}[b]{\linewidth}\centering
CV duration
\end{minipage} & \begin{minipage}[b]{\linewidth}\centering
Mean dev
\end{minipage} & \begin{minipage}[b]{\linewidth}\centering
SD dev
\end{minipage} & \begin{minipage}[b]{\linewidth}\centering
CV dev
\end{minipage} \\
\midrule\noalign{}
\endhead
\bottomrule\noalign{}
\endlastfoot
LO05 & 13 & 112.3 & 89.64 & 0.7981 & 298.4 & 306.4 & 1.027 \\
LO03 & 12 & 104.2 & 69.82 & 0.6698 & 365.6 & 306.7 & 0.8391 \\
\end{longtable}

\subsubsection{Relationship between variability and mean operative
duration}\label{relationship-between-variability-and-mean-operative-duration}

Examining how variability scales with operative duration provides
insight into whether dispersion in surgical times increases
proportionally with the length of a procedure. By analysing this
relationship at the level of procedure types, it becomes possible to
distinguish structural scaling effects from differences in planning
performance. The following figures present three complementary
perspectives: the standard deviation of operative duration, the
coefficient of variation of observed duration, and the coefficient of
variation of planning deviation. Together, they describe how both
absolute and relative variability evolve as surgeries become longer or
more complex.

The first figure shows a strong positive linear relationship between
mean operative time and its standard deviation, as illustrated in Figure
\ref{fig:sd_mean_relation}. As the average duration increases, the
absolute spread of recorded times expands almost proportionally. This
scaling effect is expected, since longer procedures allow for greater
absolute variation while maintaining a similar proportional spread. For
instance, procedures with a mean duration near one hundred minutes
typically show standard deviations around forty to fifty minutes, while
operations exceeding three hundred minutes often vary by more than one
hundred minutes. This pattern confirms that absolute variability grows
with case length without implying greater relative unpredictability.

\begin{figure}[H]

{\centering \includegraphics[width=0.7\linewidth,]{In-Depth_Analysis_Lanaken_files/figure-latex/sd mean relation-1} 

}

\caption{Standard deviation in operative time relative to mean duration \label{fig:sd_mean_relation}}\label{fig:sd mean relation}
\end{figure}

When variability is expressed relative to mean duration, a clear inverse
pattern emerges, as illustrated in Figure \ref{fig:cv_mean_relation}.
The coefficient of variation decreases gradually as procedures become
longer, indicating that proportional dispersion diminishes with case
length. Short procedures often show CV values above 0.4, while most long
operations fall below 0.3. This relationship reflects a simple scaling
property, where longer procedures display greater absolute variability
but proportionally more stability compared with their expected duration.

\begin{figure}[H]

{\centering \includegraphics[width=0.7\linewidth,]{In-Depth_Analysis_Lanaken_files/figure-latex/cv mean relation-1} 

}

\caption{Coefficient of variation (operative duration) relative to mean duration \label{fig:cv_mean_relation}}\label{fig:cv mean relation}
\end{figure}

A similar downward trend is visible when plotting the coefficient of
variation of planning deviation against mean duration, as shown in
Figure \ref{fig:cv_planning_relation}. For shorter procedures, the
relative spread in planning deviations is wide, with CV values sometimes
exceeding 3.0 to 5.0. These high coefficients occur because small
absolute deviations represent large proportional differences when the
planned duration is short. As procedure length increases, this
proportional effect diminishes, and the CV of planning deviation
stabilizes around 1.0 or lower. The negative slope therefore reflects
the same scaling principle observed for operative duration, rather than
a systematic improvement in planning accuracy.

\begin{figure}[H]

{\centering \includegraphics[width=0.7\linewidth,]{In-Depth_Analysis_Lanaken_files/figure-latex/cv planning relation-1} 

}

\caption{Coefficient of variation (planning deviation) relative to mean duration \label{fig:cv_planning_relation}}\label{fig:cv planning relation}
\end{figure}

Taken together, the three figures show that while absolute variability
in surgical time increases with case length, relative variability tends
to decrease. This pattern indicates that proportional uncertainty is
greatest for short procedures, where small timing differences represent
a large share of the total duration, and least for long surgeries, where
absolute deviations are diluted across extended operating periods.

\subsection{How large are start-time delays, and when do they
happen?}\label{how-large-are-start-time-delays-and-when-do-they-happen}

Start time deviations show a wide spread, with late starts occurring
slightly more often than early ones, as summarised in Table
\ref{tab:start_diff}. About 56.8 percent of all surgeries begin later
than scheduled, with an average delay of 32.2 minutes and a median of 21
minutes. Early starts occur in 43.2 percent of cases, with an average
lead of 33.5 minutes and a median of 20 minutes. The dispersion is
considerable: standard deviations reach 38.8 minutes for late starts and
47.3 minutes for early starts, illustrating substantial variability in
when cases begin. A few extreme outliers exceed 1,100 minutes in the
late direction, highlighting the presence of rare but very large
deviations.

These results confirm that punctuality at the start of the surgical
schedule remains a recurrent operational challenge. Even moderate delays
early in the day can cascade into later sessions, reducing scheduling
flexibility and overall room efficiency.

\begin{longtable}[]{@{}
  >{\centering\arraybackslash}p{(\linewidth - 20\tabcolsep) * \real{0.1628}}
  >{\centering\arraybackslash}p{(\linewidth - 20\tabcolsep) * \real{0.0814}}
  >{\centering\arraybackslash}p{(\linewidth - 20\tabcolsep) * \real{0.1047}}
  >{\centering\arraybackslash}p{(\linewidth - 20\tabcolsep) * \real{0.0814}}
  >{\centering\arraybackslash}p{(\linewidth - 20\tabcolsep) * \real{0.0698}}
  >{\centering\arraybackslash}p{(\linewidth - 20\tabcolsep) * \real{0.0698}}
  >{\centering\arraybackslash}p{(\linewidth - 20\tabcolsep) * \real{0.0814}}
  >{\centering\arraybackslash}p{(\linewidth - 20\tabcolsep) * \real{0.0698}}
  >{\centering\arraybackslash}p{(\linewidth - 20\tabcolsep) * \real{0.0698}}
  >{\centering\arraybackslash}p{(\linewidth - 20\tabcolsep) * \real{0.0930}}
  >{\centering\arraybackslash}p{(\linewidth - 20\tabcolsep) * \real{0.1163}}@{}}
\caption{Difference between actual and planned start time (absolute
minutes), separated by direction \label{tab:start_diff}}\tabularnewline
\toprule\noalign{}
\begin{minipage}[b]{\linewidth}\centering
direction
\end{minipage} & \begin{minipage}[b]{\linewidth}\centering
Mean
\end{minipage} & \begin{minipage}[b]{\linewidth}\centering
Median
\end{minipage} & \begin{minipage}[b]{\linewidth}\centering
SD
\end{minipage} & \begin{minipage}[b]{\linewidth}\centering
IQR
\end{minipage} & \begin{minipage}[b]{\linewidth}\centering
Min
\end{minipage} & \begin{minipage}[b]{\linewidth}\centering
Max
\end{minipage} & \begin{minipage}[b]{\linewidth}\centering
P95
\end{minipage} & \begin{minipage}[b]{\linewidth}\centering
P99
\end{minipage} & \begin{minipage}[b]{\linewidth}\centering
P99.9
\end{minipage} & \begin{minipage}[b]{\linewidth}\centering
Percent
\end{minipage} \\
\midrule\noalign{}
\endfirsthead
\toprule\noalign{}
\begin{minipage}[b]{\linewidth}\centering
direction
\end{minipage} & \begin{minipage}[b]{\linewidth}\centering
Mean
\end{minipage} & \begin{minipage}[b]{\linewidth}\centering
Median
\end{minipage} & \begin{minipage}[b]{\linewidth}\centering
SD
\end{minipage} & \begin{minipage}[b]{\linewidth}\centering
IQR
\end{minipage} & \begin{minipage}[b]{\linewidth}\centering
Min
\end{minipage} & \begin{minipage}[b]{\linewidth}\centering
Max
\end{minipage} & \begin{minipage}[b]{\linewidth}\centering
P95
\end{minipage} & \begin{minipage}[b]{\linewidth}\centering
P99
\end{minipage} & \begin{minipage}[b]{\linewidth}\centering
P99.9
\end{minipage} & \begin{minipage}[b]{\linewidth}\centering
Percent
\end{minipage} \\
\midrule\noalign{}
\endhead
\bottomrule\noalign{}
\endlastfoot
Late start & 32.2 & 21 & 38.8 & 32 & 1 & 1119 & 99 & 186 & 338.2 &
56.8\% \\
Early start & 33.5 & 20 & 47.3 & 34 & 0 & 822 & 113 & 225 & 462 &
43.2\% \\
\end{longtable}

The distribution of start-time differences is centered around zero, as
shown in Figure \ref{fig:start_diff_hist}, with a clear peak at small
deviations, indicating that most surgeries start close to their planned
time. However, the shape of the histogram is slightly skewed toward
positive values, showing that late starts are more frequent than early
ones. Most cases begin within approximately half an hour of schedule,
but the gradual widening of the distribution toward the right reveals
that a smaller subset experiences significant delays.

The presence of a long right-hand tail confirms that occasionally,
extreme delays still occur. These outliers have a limited frequency but
can affect overall operating room throughput and schedule stability. The
symmetrical appearance around the center, combined with a mild skew,
suggests that deviations are not random but follow a consistent pattern
of minor delay accumulation across the day.

\begin{figure}[H]

{\centering \includegraphics[width=0.7\linewidth,]{In-Depth_Analysis_Lanaken_files/figure-latex/start diff hist-1} 

}

\caption{Start time deviations (grouped per 10 minutes) \label{fig:start_diff_hist}}\label{fig:start diff hist}
\end{figure}

Average start-time delays show moderate variation across weekdays, as
illustrated in Figure \ref{fig:start_delay_weekday}, ranging from minus
1.8 to 7.4 minutes during regular working days. Monday and Tuesday have
mean delays of 7.4 and 6.1 minutes, while Wednesday reaches 6.8 minutes.
Thursday and Friday follow with minus 0.7 and minus 1.8 minutes,
respectively. These differences indicate that start-time punctuality is
relatively consistent throughout the first three days of the week, with
fluctuations on Thursday and Friday that are likely to reflect
structural scheduling issues.

\begin{figure}[H]

{\centering \includegraphics[width=0.7\linewidth,]{In-Depth_Analysis_Lanaken_files/figure-latex/start delay weekday-1} 

}

\caption{Mean start delay by weekday \label{fig:start_delay_weekday}}\label{fig:start delay weekday}
\end{figure}

Late starts occur more frequently than early starts on every day of the
week, as summarized in Table \ref{tab:start_diff_weekday}. Across
regular weekdays, about three fifths of surgeries begin later than
planned, with mean delays between 28.6 and 35.4 minutes and median
delays around 18 to 23 minutes. Early starts are less common,
representing roughly 37.6 to 46.7\% of all cases, and typically begin
between 32.1 and 40 minutes before the scheduled time. These small
variations between weekdays indicate that overall punctuality during the
workweek is relatively stable, with no clear day standing out as an
outlier.

\begin{longtable}[]{@{}
  >{\centering\arraybackslash}p{(\linewidth - 14\tabcolsep) * \real{0.0917}}
  >{\centering\arraybackslash}p{(\linewidth - 14\tabcolsep) * \real{0.0734}}
  >{\centering\arraybackslash}p{(\linewidth - 14\tabcolsep) * \real{0.1651}}
  >{\centering\arraybackslash}p{(\linewidth - 14\tabcolsep) * \real{0.1193}}
  >{\centering\arraybackslash}p{(\linewidth - 14\tabcolsep) * \real{0.1376}}
  >{\centering\arraybackslash}p{(\linewidth - 14\tabcolsep) * \real{0.1743}}
  >{\centering\arraybackslash}p{(\linewidth - 14\tabcolsep) * \real{0.1101}}
  >{\centering\arraybackslash}p{(\linewidth - 14\tabcolsep) * \real{0.1284}}@{}}
\caption{Start-time deviation patterns per weekday (minutes).
\label{tab:start_diff_weekday}}\tabularnewline
\toprule\noalign{}
\begin{minipage}[b]{\linewidth}\centering
Weekday
\end{minipage} & \begin{minipage}[b]{\linewidth}\centering
n
\end{minipage} & \begin{minipage}[b]{\linewidth}\centering
Late starts (\%)
\end{minipage} & \begin{minipage}[b]{\linewidth}\centering
Mean delay
\end{minipage} & \begin{minipage}[b]{\linewidth}\centering
Median delay
\end{minipage} & \begin{minipage}[b]{\linewidth}\centering
Early starts (\%)
\end{minipage} & \begin{minipage}[b]{\linewidth}\centering
Mean lead
\end{minipage} & \begin{minipage}[b]{\linewidth}\centering
Median lead
\end{minipage} \\
\midrule\noalign{}
\endfirsthead
\toprule\noalign{}
\begin{minipage}[b]{\linewidth}\centering
Weekday
\end{minipage} & \begin{minipage}[b]{\linewidth}\centering
n
\end{minipage} & \begin{minipage}[b]{\linewidth}\centering
Late starts (\%)
\end{minipage} & \begin{minipage}[b]{\linewidth}\centering
Mean delay
\end{minipage} & \begin{minipage}[b]{\linewidth}\centering
Median delay
\end{minipage} & \begin{minipage}[b]{\linewidth}\centering
Early starts (\%)
\end{minipage} & \begin{minipage}[b]{\linewidth}\centering
Mean lead
\end{minipage} & \begin{minipage}[b]{\linewidth}\centering
Median lead
\end{minipage} \\
\midrule\noalign{}
\endhead
\bottomrule\noalign{}
\endlastfoot
Ma & 14710 & 59.7 & 32.8 & 21 & 37.9 & -32.1 & -19 \\
Di & 14902 & 60.1 & 35.4 & 23 & 37.8 & -40 & -22 \\
Wo & 13742 & 60 & 32.1 & 21 & 37.6 & -33.2 & -20 \\
Do & 13233 & 51.7 & 28.6 & 18 & 45.9 & -33.8 & -22 \\
Vr & 12808 & 51.2 & 30.7 & 20 & 46.7 & -37.6 & -23 \\
\end{longtable}

Start-time reliability varies widely between operating rooms, as shown
in Table \ref{tab:start_diff_or}. The proportion of late starts ranges
from 47.5\% to 66\%, indicating that punctuality is evenly distributed
across the OR complex. Some rooms, such as LO04 and LP01, stand out with
high shares of late starts, 66\% and 63.5\%, respectively. In LO03, the
mean delay among late cases reaches 50.5 minutes, driven by a small
number of outliers that strongly influence the average.

At the opposite end, rooms such as LO01 and LP02 display a balanced
distribution, where early starts occur roughly as often as late ones.
These contrasts suggest that certain rooms operate under more stable and
predictable timing conditions, while others experience recurring or
structural delays. Overall, start-time deviations appear concentrated in
a limited set of operating rooms, where occasional long postponements
account for much of the overall variability in punctuality.

\begin{longtable}[]{@{}
  >{\centering\arraybackslash}p{(\linewidth - 14\tabcolsep) * \real{0.1000}}
  >{\centering\arraybackslash}p{(\linewidth - 14\tabcolsep) * \real{0.0727}}
  >{\centering\arraybackslash}p{(\linewidth - 14\tabcolsep) * \real{0.1636}}
  >{\centering\arraybackslash}p{(\linewidth - 14\tabcolsep) * \real{0.1182}}
  >{\centering\arraybackslash}p{(\linewidth - 14\tabcolsep) * \real{0.1364}}
  >{\centering\arraybackslash}p{(\linewidth - 14\tabcolsep) * \real{0.1727}}
  >{\centering\arraybackslash}p{(\linewidth - 14\tabcolsep) * \real{0.1091}}
  >{\centering\arraybackslash}p{(\linewidth - 14\tabcolsep) * \real{0.1273}}@{}}
\caption{Start-time deviation patterns per operating room (minutes).
\label{tab:start_diff_or}}\tabularnewline
\toprule\noalign{}
\begin{minipage}[b]{\linewidth}\centering
ActualOR
\end{minipage} & \begin{minipage}[b]{\linewidth}\centering
n
\end{minipage} & \begin{minipage}[b]{\linewidth}\centering
Late starts (\%)
\end{minipage} & \begin{minipage}[b]{\linewidth}\centering
Mean delay
\end{minipage} & \begin{minipage}[b]{\linewidth}\centering
Median delay
\end{minipage} & \begin{minipage}[b]{\linewidth}\centering
Early starts (\%)
\end{minipage} & \begin{minipage}[b]{\linewidth}\centering
Mean lead
\end{minipage} & \begin{minipage}[b]{\linewidth}\centering
Median lead
\end{minipage} \\
\midrule\noalign{}
\endfirsthead
\toprule\noalign{}
\begin{minipage}[b]{\linewidth}\centering
ActualOR
\end{minipage} & \begin{minipage}[b]{\linewidth}\centering
n
\end{minipage} & \begin{minipage}[b]{\linewidth}\centering
Late starts (\%)
\end{minipage} & \begin{minipage}[b]{\linewidth}\centering
Mean delay
\end{minipage} & \begin{minipage}[b]{\linewidth}\centering
Median delay
\end{minipage} & \begin{minipage}[b]{\linewidth}\centering
Early starts (\%)
\end{minipage} & \begin{minipage}[b]{\linewidth}\centering
Mean lead
\end{minipage} & \begin{minipage}[b]{\linewidth}\centering
Median lead
\end{minipage} \\
\midrule\noalign{}
\endhead
\bottomrule\noalign{}
\endlastfoot
LO04 & 6009 & 66 & 43.2 & 26 & 32.1 & -52.2 & -30 \\
LP01 & 20189 & 63.5 & 24.9 & 20 & 34.8 & -24.7 & -16 \\
LO03 & 4196 & 59.6 & 50.5 & 31 & 38 & -53.2 & -31 \\
LO02 & 9974 & 57.1 & 44.7 & 29 & 38.4 & -56.3 & -33 \\
LO05 & 5206 & 56.9 & 41.3 & 23 & 40.7 & -45.3 & -26 \\
LO01 & 5202 & 49.9 & 33.4 & 21 & 47.6 & -35.3 & -20 \\
LP02 & 18619 & 47.5 & 21.1 & 16 & 50.9 & -26.2 & -20 \\
\end{longtable}

\subsection{Which procedures deviate most from
plan?}\label{which-procedures-deviate-most-from-plan}

Procedures differ noticeably in how accurately their duration is
predicted, as shown in Table \ref{tab:duration_dev_proc}. The twenty
procedure types with the largest absolute deviations display mean
differences ranging from about 21 to nearly 39 minutes, indicating that
for several procedure types the planned duration can diverge
meaningfully from actual performance.

The largest deviations appear in VARICES VSM and LASER -- MK, where mean
absolute differences exceed 38 minutes. These procedures stand out as
the most difficult to estimate accurately, likely because their
execution times vary with patient-specific factors or case complexity.
Several other procedures with moderate case volumes, such as SCHOUDEER,
BREDE EXCISIE MALIGNE, and EXPLORATIEVE LAPARASCOPIE, show deviations in
the range of 26 to 29 minutes, reflecting a mix of predictable structure
with occasional outliers.

More frequently performed procedures, including VARICES, VARICES LASER,
and PEES HAND (VINGER), show smaller deviations around 21 minutes,
suggesting that repetition and routine support more accurate planning.
The pattern therefore aligns with an experience-based effect: procedures
performed often tend to have tighter planning accuracy, while less
frequent or clinically variable procedures show greater dispersion.

Across the listed procedures, the share of longer-than-planned cases
exceeds the share of shorter ones for most items. This indicates a
systematic tendency toward underestimation, although the magnitude
varies between procedure types.

Overall, these results show that planning accuracy at Campus Lanaken
depends strongly on procedure type. Common, routine interventions tend
to follow their planned durations closely, while specialized or more
variable procedures exhibit wider deviations even when the average case
length is not particularly long.

\begin{longtable}[]{@{}
  >{\centering\arraybackslash}p{(\linewidth - 12\tabcolsep) * \real{0.2619}}
  >{\centering\arraybackslash}p{(\linewidth - 12\tabcolsep) * \real{0.0556}}
  >{\centering\arraybackslash}p{(\linewidth - 12\tabcolsep) * \real{0.1667}}
  >{\centering\arraybackslash}p{(\linewidth - 12\tabcolsep) * \real{0.1032}}
  >{\centering\arraybackslash}p{(\linewidth - 12\tabcolsep) * \real{0.1508}}
  >{\centering\arraybackslash}p{(\linewidth - 12\tabcolsep) * \real{0.1111}}
  >{\centering\arraybackslash}p{(\linewidth - 12\tabcolsep) * \real{0.1508}}@{}}
\caption{Top 20 procedures by average absolute duration deviation
(Lanaken). \label{tab:duration_dev_proc}}\tabularnewline
\toprule\noalign{}
\begin{minipage}[b]{\linewidth}\centering
ProcedureName
\end{minipage} & \begin{minipage}[b]{\linewidth}\centering
n
\end{minipage} & \begin{minipage}[b]{\linewidth}\centering
Mean abs dev (min)
\end{minipage} & \begin{minipage}[b]{\linewidth}\centering
Longer (\%)
\end{minipage} & \begin{minipage}[b]{\linewidth}\centering
Mean + dev (min)
\end{minipage} & \begin{minipage}[b]{\linewidth}\centering
Shorter (\%)
\end{minipage} & \begin{minipage}[b]{\linewidth}\centering
Mean - dev (min)
\end{minipage} \\
\midrule\noalign{}
\endfirsthead
\toprule\noalign{}
\begin{minipage}[b]{\linewidth}\centering
ProcedureName
\end{minipage} & \begin{minipage}[b]{\linewidth}\centering
n
\end{minipage} & \begin{minipage}[b]{\linewidth}\centering
Mean abs dev (min)
\end{minipage} & \begin{minipage}[b]{\linewidth}\centering
Longer (\%)
\end{minipage} & \begin{minipage}[b]{\linewidth}\centering
Mean + dev (min)
\end{minipage} & \begin{minipage}[b]{\linewidth}\centering
Shorter (\%)
\end{minipage} & \begin{minipage}[b]{\linewidth}\centering
Mean - dev (min)
\end{minipage} \\
\midrule\noalign{}
\endhead
\bottomrule\noalign{}
\endlastfoot
VARICES VSM LASER - MK & 49 & 38.6 & 73.5 & 34.4 & 26.5 & -50.3 \\
SCHOUDER ROTATOR CUFF ARTROSCOPISCH & 165 & 28.7 & 46.1 & 32.8 & 52.7 &
-25.9 \\
CHIR.IMPLNT.ACT.IMPL & 66 & 28 & 33.3 & 12.2 & 66.7 & -36 \\
EXPLORATIEVE LAPARASCOPIE & 47 & 26.8 & 12.8 & 11.3 & 87.2 & -29 \\
BREDE EXCISIE MALIGNE MELANOOM MET SENT KLIERBIOPS & 44 & 26.5 & 63.6 &
34.8 & 34.1 & -12.7 \\
CLAVICULA FRAKTUUR AO KLEIN FRAGMENT & 47 & 26.4 & 53.2 & 27.7 & 44.7 &
-26.2 \\
Ganglion van Gasser : radiofrequente R & 88 & 26.1 & 1.1 & 4 & 98.9 &
-26.3 \\
HEELKUNDIGE CORRECTIE VAN GYNAECOMASTIE & 58 & 25.8 & 62.1 & 29.5 & 36.2
& -20.7 \\
REDUCERENDE MAMMAPLASTIEK & 73 & 24.6 & 46.6 & 29 & 50.7 & -21.8 \\
LITTEKENCORRECTIE & 65 & 24.1 & 53.8 & 31.2 & 41.5 & -17.5 \\
ABDOMENPLASTIE & 44 & 23.4 & 59.1 & 25.7 & 38.6 & -21.2 \\
HAND DUPYTREN & 75 & 22.9 & 33.3 & 16.4 & 64 & -27.2 \\
DESCEMET'S STRIPPING AUTOMATED ENDOTHELIAL KERTOPLASTY -
ENDOTHEELTRANSPLANTATIE & 34 & 22.3 & 47.1 & 27.8 & 50 & -18.4 \\
epidurale stimulatie:diagnostisch lumbaal enkele elektrode & 41 & 22.1 &
63.4 & 26.1 & 36.6 & -15.2 \\
Ganglion sfenopalatinum : Radiofrequente R & 78 & 22 & 0 & NA & 98.7 &
-22.3 \\
PEES HAND (VINGER) VOET HECHTING & 46 & 21.7 & 26.1 & 26.6 & 71.7 &
-20.5 \\
VARICES & 689 & 21.3 & 52.5 & 25.2 & 46 & -17.5 \\
VARICES LASER & 1351 & 21.3 & 49.3 & 24.8 & 49.4 & -18.4 \\
EXC MALIGNE HUIDLETSEL MT RECONST HUIDGR HFD/GEZIC & 45 & 21.2 & 48.9 &
23.8 & 51.1 & -18.7 \\
epidurale stimulatie: batterij implantatie & 54 & 21.1 & 64.8 & 25 &
33.3 & -14.9 \\
\end{longtable}

The scatter plot in Figure \ref{fig:volume_deviation} shows an inverse
relationship between procedure volume and average deviation. Procedures
that occur more frequently tend to have lower mean absolute deviations,
while rare procedures show much greater variability in duration. This
pattern is visible in the downward trend of the points as procedure
volume increases.

High-volume procedures, which appear toward the right of the chart,
cluster at lower deviation values, indicating more stable and
predictable planning performance. In contrast, low-volume procedures on
the left display a much wider range of deviations that can exceed 25
minutes. The relationship suggests that familiarity and repetition
contribute to more accurate duration estimates, whereas infrequent
procedures are harder to predict due to limited historical data.
Overall, the figure highlights that case volume is an important factor
in explaining the variability of planning accuracy across procedure
types.

\begin{figure}[H]

{\centering \includegraphics[width=0.7\linewidth,]{In-Depth_Analysis_Lanaken_files/figure-latex/volume_deviation-1} 

}

\caption{Relationship between procedure volume and mean deviation \label{fig:volume_deviation}}\label{fig:volume_deviation}
\end{figure}

The distribution of duration deviations across surgeons reveals notable
variation in planning accuracy, as shown in Table
\ref{tab:surgeon_deviation}. Among those with at least fifty recorded
procedures, the average absolute deviation ranges from 15.7 to 25
minutes, while the average relative deviation spans 20.2\% to 28.5\%. A
few surgeons show comparatively lower mean deviations below 19 minutes,
whereas most fall between 19 and 24 minutes, indicating that planning
performance differs across individuals but remains within a moderate
range for the majority.

No clear relationship appears between the number of recorded cases and
the size of the deviation, suggesting that factors beyond case volume
influence the observed variation. Overall, the results show that
duration deviations are not uniformly distributed across surgeons,
emphasizing that planning accuracy can vary meaningfully even within a
shared operational environment.

\begin{longtable}[]{@{}lccc@{}}
\caption{Top 15 surgeons by average absolute and relative duration
deviation (pseudonymized). \label{tab:surgeon_deviation}}\tabularnewline
\toprule\noalign{}
Surgeon & n & Mean absolute dev (min) & Mean relative dev (\%) \\
\midrule\noalign{}
\endfirsthead
\toprule\noalign{}
Surgeon & n & Mean absolute dev (min) & Mean relative dev (\%) \\
\midrule\noalign{}
\endhead
\bottomrule\noalign{}
\endlastfoot
Surgeon 1 & 353 & 25.0 & 25.9 \\
Surgeon 2 & 223 & 24.5 & 25.1 \\
Surgeon 3 & 597 & 24.2 & 25.4 \\
Surgeon 4 & 728 & 23.9 & 24.5 \\
Surgeon 5 & 151 & 23.3 & 26.8 \\
Surgeon 6 & 361 & 21.8 & 28.5 \\
Surgeon 7 & 73 & 21.6 & 28.0 \\
Surgeon 8 & 202 & 20.5 & 21.8 \\
Surgeon 9 & 603 & 19.0 & 23.2 \\
Surgeon 10 & 128 & 18.7 & 22.9 \\
Surgeon 11 & 270 & 18.6 & 20.9 \\
Surgeon 12 & 376 & 18.4 & 20.2 \\
Surgeon 13 & 361 & 16.5 & 24.2 \\
Surgeon 14 & 132 & 16.4 & 21.6 \\
Surgeon 15 & 297 & 15.7 & 25.3 \\
\end{longtable}

\newpage

\section{Do surgeries stay within regular working
hours?}\label{do-surgeries-stay-within-regular-working-hours}

\subsection{How often do cases extend beyond their assigned time
window?}\label{how-often-do-cases-extend-beyond-their-assigned-time-window}

Only a very small share of surgeries at Lanaken extends beyond the
scheduled shift end, as shown in Table \ref{tab:overtime_summary}. Out
of nearly seventy thousand procedures, 0.5 percent continue into
after-hours, indicating that overtime is an exceptional rather than
routine event in this campus. When it does occur, overtime remains
limited: the mean extension is 15.7 minutes, and the median is 10
minutes, showing that most overruns are brief and operationally
contained. A smaller subset of cases lasts longer, with the 95th
percentile reaching 48.5 minutes, illustrating that occasional outliers
contribute disproportionately to the total overtime minutes.

These results highlight a scheduling system that operates with very high
temporal adherence. The extremely low percentage of after-hours cases
suggests that the daily program rarely spills over into the next shift,
and when it does, the additional time required is usually modest. The
distribution nevertheless shows that a few atypical procedures or
unexpected delays can still push individual cases into the after-hours
window, but such events remain rare within the overall workflow at
Lanaken.

\begin{longtable}[]{@{}
  >{\centering\arraybackslash}p{(\linewidth - 10\tabcolsep) * \real{0.0833}}
  >{\centering\arraybackslash}p{(\linewidth - 10\tabcolsep) * \real{0.2708}}
  >{\centering\arraybackslash}p{(\linewidth - 10\tabcolsep) * \real{0.1354}}
  >{\centering\arraybackslash}p{(\linewidth - 10\tabcolsep) * \real{0.1667}}
  >{\centering\arraybackslash}p{(\linewidth - 10\tabcolsep) * \real{0.1875}}
  >{\centering\arraybackslash}p{(\linewidth - 10\tabcolsep) * \real{0.1562}}@{}}
\caption{Overall share and duration of after-hours surgeries (Lanaken).
\label{tab:overtime_summary}}\tabularnewline
\toprule\noalign{}
\begin{minipage}[b]{\linewidth}\centering
Total
\end{minipage} & \begin{minipage}[b]{\linewidth}\centering
Surgeries with overtime
\end{minipage} & \begin{minipage}[b]{\linewidth}\centering
Percentage
\end{minipage} & \begin{minipage}[b]{\linewidth}\centering
Mean Overtime
\end{minipage} & \begin{minipage}[b]{\linewidth}\centering
Median Overtime
\end{minipage} & \begin{minipage}[b]{\linewidth}\centering
P95 Overtime
\end{minipage} \\
\midrule\noalign{}
\endfirsthead
\toprule\noalign{}
\begin{minipage}[b]{\linewidth}\centering
Total
\end{minipage} & \begin{minipage}[b]{\linewidth}\centering
Surgeries with overtime
\end{minipage} & \begin{minipage}[b]{\linewidth}\centering
Percentage
\end{minipage} & \begin{minipage}[b]{\linewidth}\centering
Mean Overtime
\end{minipage} & \begin{minipage}[b]{\linewidth}\centering
Median Overtime
\end{minipage} & \begin{minipage}[b]{\linewidth}\centering
P95 Overtime
\end{minipage} \\
\midrule\noalign{}
\endhead
\bottomrule\noalign{}
\endlastfoot
69395 & 371 & 0.5 & 15.7 & 10 & 48.5 \\
\end{longtable}

The distribution of overtime duration, shown in Figure
\ref{fig:overtime_dist}, is strongly right-skewed. Most extended cases
end within the first 30 minutes beyond the scheduled shift, and the
highest frequency occurs between 0 and 20 minutes of overtime. Beyond
that, the number of cases declines steadily, with only a small fraction
exceeding 60 minutes. A few extreme cases surpass 100 minutes, but these
are rare.

This pattern confirms that while overtime is a recurring operational
feature, it is typically limited in duration. The long right tail
illustrates that a small group of prolonged procedures accounts for a
disproportionate share of total overtime minutes. These outliers can
heavily affect staff workload and room turnover late in the day, but
they do not represent the typical case. Overall, overtime at Lanaken
mainly consists of brief extensions of planned activity rather than
frequent, sustained night work.

\begin{figure}[H]

{\centering \includegraphics[width=0.7\linewidth,]{In-Depth_Analysis_Lanaken_files/figure-latex/overtime_dist-1} 

}

\caption{Distribution of overtime duration for after-hours surgeries \label{fig:overtime_dist}}\label{fig:overtime_dist}
\end{figure}

The share of surgeries extending beyond scheduled shift hours remains
steady throughout the regular workweek, as shown in Figure
\ref{fig:afterhours_weekday}, fluctuating between 0.3 and 0.8 percent
from Monday to Friday. This indicates that weekday operations are
largely contained within their planned time frames, with only a small
proportion of cases requiring additional time beyond shift boundaries.

\begin{figure}[H]

{\centering \includegraphics[width=0.7\linewidth,]{In-Depth_Analysis_Lanaken_files/figure-latex/afterhours_weekday-1} 

}

\caption{Share of after-hours surgeries by weekday \label{fig:afterhours_weekday}}\label{fig:afterhours_weekday}
\end{figure}

The proportion of surgeries extending beyond scheduled shift hours
remains extremely low and shows no structural change over time, as
illustrated in Figure \ref{fig:afterhours_trend}. In 2022 and 2023,
about 0.5 percent of surgeries at Lanaken exceeded their scheduled
shift. This share increases slightly to 0.6 percent in 2024 before
returning to 0.5 percent in the partial dataset for 2025.

These values represent the percentage of all procedures performed in
each year that continued past the official end of their assigned shift.
The near-identical year-to-year percentages demonstrate that after-hours
activity at Lanaken is highly stable, with only minimal variation across
the observed period.

Because overtime is so rare, these small fluctuations are unlikely to
signal operational changes. Instead, they probably reflect routine
variability in low-frequency events rather than shifts in throughput,
planning accuracy, or case mix. Overall, after-hours work at Lanaken is
a consistently marginal component of surgical activity, indicating very
strong adherence to scheduled shift boundaries and efficient daily
workflow management.

\begin{figure}[H]

{\centering \includegraphics[width=0.7\linewidth,]{In-Depth_Analysis_Lanaken_files/figure-latex/afterhours_trend-1} 

}

\caption{Evolution of after-hours activity over time \label{fig:afterhours_trend}}\label{fig:afterhours_trend}
\end{figure}

The proportion of surgeries extending beyond scheduled shift hours
differs markedly across operating rooms, as summarized in Table
\ref{tab:overtime_or}. Most rooms show overtime rates between 0.6 and
1.2 percent, while one stands out with slightly higher shares. LO03
records the highest proportion of overtime cases at 1.7 percent, with a
mean duration of 18.9 minutes and a 95th percentile of 51.8 minutes,
indicating frequent and occasionally prolonged extensions. Such high
figures likely reflect the type of procedures performed in this room,
which are often complex and less predictable in length.

Several other rooms, including LO05, LO01, and LO04, show overtime in
roughly 1.1 to 1.2 percent of cases, with mean durations between 15.2
and 18 minutes. These values suggest that while most surgeries in these
rooms finish within their allocated shift, a small share still extends
well into the next operational block.

Overall, the data indicate that overtime at Lanaken is concentrated in a
limited number of high-volume or complex operating rooms. While many
rooms complete their scheduled activity within the planned shift most of
the time, a small subset still experiences overtime in 0.1 to 1.7
percent of surgeries, showing that after-shift extensions remain a
recurring part of daily operations rather than isolated exceptions.

\begin{longtable}[]{@{}
  >{\centering\arraybackslash}p{(\linewidth - 10\tabcolsep) * \real{0.1294}}
  >{\centering\arraybackslash}p{(\linewidth - 10\tabcolsep) * \real{0.0941}}
  >{\centering\arraybackslash}p{(\linewidth - 10\tabcolsep) * \real{0.2941}}
  >{\centering\arraybackslash}p{(\linewidth - 10\tabcolsep) * \real{0.1882}}
  >{\centering\arraybackslash}p{(\linewidth - 10\tabcolsep) * \real{0.2118}}
  >{\centering\arraybackslash}p{(\linewidth - 10\tabcolsep) * \real{0.0824}}@{}}
\caption{Operating rooms by after-hours activity (Lanaken).
\label{tab:overtime_or}}\tabularnewline
\toprule\noalign{}
\begin{minipage}[b]{\linewidth}\centering
ActualOR
\end{minipage} & \begin{minipage}[b]{\linewidth}\centering
Cases
\end{minipage} & \begin{minipage}[b]{\linewidth}\centering
Percentage of overtime surgeries
\end{minipage} & \begin{minipage}[b]{\linewidth}\centering
Mean Overtime
\end{minipage} & \begin{minipage}[b]{\linewidth}\centering
Median Overtime
\end{minipage} & \begin{minipage}[b]{\linewidth}\centering
P95
\end{minipage} \\
\midrule\noalign{}
\endfirsthead
\toprule\noalign{}
\begin{minipage}[b]{\linewidth}\centering
ActualOR
\end{minipage} & \begin{minipage}[b]{\linewidth}\centering
Cases
\end{minipage} & \begin{minipage}[b]{\linewidth}\centering
Percentage of overtime surgeries
\end{minipage} & \begin{minipage}[b]{\linewidth}\centering
Mean Overtime
\end{minipage} & \begin{minipage}[b]{\linewidth}\centering
Median Overtime
\end{minipage} & \begin{minipage}[b]{\linewidth}\centering
P95
\end{minipage} \\
\midrule\noalign{}
\endhead
\bottomrule\noalign{}
\endlastfoot
LO03 & 4196 & 1.7 & 18.9 & 14 & 51.8 \\
LO05 & 5206 & 1.2 & 18 & 15 & 46 \\
LO01 & 5202 & 1.1 & 20.7 & 15 & 63.8 \\
LO04 & 6009 & 1.1 & 15.2 & 9 & 37.6 \\
LO02 & 9974 & 0.6 & 12 & 9 & 38.7 \\
LP01 & 20189 & 0.1 & 9.1 & 8 & 19.8 \\
LP02 & 18619 & 0.1 & 7.7 & 5 & 19.6 \\
\end{longtable}

\subsection{When does overtime occur, and how heavy is
it?}\label{when-does-overtime-occur-and-how-heavy-is-it}

The distribution of overtime surgeries, shown in Figure
\ref{fig:overtime_timing}, displays a clear concentration of cases
ending shortly after 16:30. This peak primarily represents procedures
that started during the regular day shift and finished later than
planned. The number of surgeries declines rapidly as the evening
progresses, indicating that most cases extending beyond their scheduled
shift end within the first one to two hours afterward.

After 22:00, only a smaller number of surgeries continue to appear,
suggesting that operations ending during the late evening or night are
relatively infrequent. Very few procedures extend past midnight, and
only isolated cases finish after 07:00.

Overall, the pattern demonstrates that overtime activity is dominated by
operations that conclude soon after the end of the regular working day,
while surgeries finishing later in the night represent a much smaller
portion of total activity.

\begin{figure}[H]

{\centering \includegraphics[width=0.7\linewidth,]{In-Depth_Analysis_Lanaken_files/figure-latex/overtime_timing-1} 

}

\caption{Timing of overtime surgeries \label{fig:overtime_timing}}\label{fig:overtime_timing}
\end{figure}

Average overtime durations are highly consistent across weekdays, as
illustrated in Figure \ref{fig:overtime_weekday}, with means ranging
from 13.1 to 18.9 minutes from Monday to Friday. The median per weekday
is clearly lower, between 7.5 and 12 minutes, which indicates a
right-skewed distribution where a small number of long cases raises the
mean.

\begin{figure}[H]

{\centering \includegraphics[width=0.7\linewidth,]{In-Depth_Analysis_Lanaken_files/figure-latex/overtime_weekday-1} 

}

\caption{Average and median overtime duration by weekday \label{fig:overtime_weekday}}\label{fig:overtime_weekday}
\end{figure}

Average and median overtime durations per surgery show a largely stable
pattern over the observed years, as illustrated in Figure
\ref{fig:overtime_year}. The mean overtime remains close to 14--15
minutes, ranging from 14.1 minutes in 2022 to a temporary rise of 19.2
minutes in 2024, before returning to 14.4 minutes in 2025. The median
follows a similar structure at lower levels, fluctuating between 9 and
13.5 minutes.

The consistent gap between the mean and median indicates a right-skewed
distribution, where most overtime periods are short but a small number
of longer extensions pull the average upward. In practical terms, the
typical overtime duration (median) lies around 10 minutes, while the
mean reflects the influence of occasional longer cases.

The visible increase in 2024 represents a short-term deviation rather
than a structural shift. Across the full period, the overall level of
overtime duration remains remarkably stable, with only modest
year-to-year variation. As with other campuses, the persistence of a
mean exceeding the median underscores that a small minority of prolonged
cases accounts for a disproportionate share of total overtime minutes,
even in an environment where overtime is rare.

\begin{figure}[H]

{\centering \includegraphics[width=0.7\linewidth,]{In-Depth_Analysis_Lanaken_files/figure-latex/overtime_year-1} 

}

\caption{Average and median overtime per surgery by year \label{fig:overtime_year}}\label{fig:overtime_year}
\end{figure}

The distribution of overtime differs noticeably between shift types, as
shown in Figure \ref{fig:overtime_shift}. The day shift (08:00--16:30)
displays a wide spread, with most cases extending only slightly beyond
the scheduled end but a small number continuing for several hours. This
pattern explains why the day shift accounts for most overtime minutes
overall.

Overall, the distributions illustrate that overtime is most variable
during the day, when activity levels are highest, and more contained
during the later shifts. The pronounced right tails in each group
confirm that a limited number of long-running cases account for a
substantial share of total overtime minutes.

\begin{figure}[H]

{\centering \includegraphics[width=0.7\linewidth,]{In-Depth_Analysis_Lanaken_files/figure-latex/overtime_shift-1} 

}

\caption{Relative overtime per shift type \label{fig:overtime_shift}}\label{fig:overtime_shift}
\end{figure}

\subsection{How does Lanaken compare with Genk, Cathlab, and
Maaseik?}\label{how-does-lanaken-compare-with-genk-cathlab-and-maaseik}

The comparison between campuses reveals pronounced differences in
after-hours surgical activity. As shown in Table
\ref{tab:afterhours_campus}, Genk performs the largest share of overtime
surgeries, with 8.4 percent of all cases extending beyond 16:30. Cathlab
follows with 5.8 percent, while Maaseik records a more moderate 2.8
percent. Lanaken, by contrast, shows minimal after-hours activity at 0.5
percent, reflecting its lower surgical volume and the limited number of
complex or time-intensive cases performed there.

Overtime duration metrics further underline these contrasts. Genk shows
the highest average and upper-percentile values, with an average
overtime of 59 minutes and a ninety-fifth percentile of 193.8 minutes,
indicating that prolonged extensions are not uncommon at the main
campus. Maaseik and Cathlab display intermediate levels, while Lanaken
again shows the shortest durations, with an average overtime of 15.7
minutes and a ninety-fifth percentile of 48.5 minutes.

Taken together, these figures show that after-hours surgical activity is
heavily concentrated at the larger campuses, particularly Genk, while
Lanaken contributes only marginally to extended operating room use
within the ZOL network.

\begin{longtable}[]{@{}
  >{\centering\arraybackslash}p{(\linewidth - 12\tabcolsep) * \real{0.0820}}
  >{\centering\arraybackslash}p{(\linewidth - 12\tabcolsep) * \real{0.0656}}
  >{\centering\arraybackslash}p{(\linewidth - 12\tabcolsep) * \real{0.1803}}
  >{\centering\arraybackslash}p{(\linewidth - 12\tabcolsep) * \real{0.2459}}
  >{\centering\arraybackslash}p{(\linewidth - 12\tabcolsep) * \real{0.1557}}
  >{\centering\arraybackslash}p{(\linewidth - 12\tabcolsep) * \real{0.1475}}
  >{\centering\arraybackslash}p{(\linewidth - 12\tabcolsep) * \real{0.1230}}@{}}
\caption{Comparison of after-hours activity between campuses.
\label{tab:afterhours_campus}}\tabularnewline
\toprule\noalign{}
\begin{minipage}[b]{\linewidth}\centering
Campus
\end{minipage} & \begin{minipage}[b]{\linewidth}\centering
Total
\end{minipage} & \begin{minipage}[b]{\linewidth}\centering
AfterHour surgeries
\end{minipage} & \begin{minipage}[b]{\linewidth}\centering
Share of overtime surgeries
\end{minipage} & \begin{minipage}[b]{\linewidth}\centering
Average Overtime
\end{minipage} & \begin{minipage}[b]{\linewidth}\centering
Median Overtime
\end{minipage} & \begin{minipage}[b]{\linewidth}\centering
P95 Overtime
\end{minipage} \\
\midrule\noalign{}
\endfirsthead
\toprule\noalign{}
\begin{minipage}[b]{\linewidth}\centering
Campus
\end{minipage} & \begin{minipage}[b]{\linewidth}\centering
Total
\end{minipage} & \begin{minipage}[b]{\linewidth}\centering
AfterHour surgeries
\end{minipage} & \begin{minipage}[b]{\linewidth}\centering
Share of overtime surgeries
\end{minipage} & \begin{minipage}[b]{\linewidth}\centering
Average Overtime
\end{minipage} & \begin{minipage}[b]{\linewidth}\centering
Median Overtime
\end{minipage} & \begin{minipage}[b]{\linewidth}\centering
P95 Overtime
\end{minipage} \\
\midrule\noalign{}
\endhead
\bottomrule\noalign{}
\endlastfoot
Genk & 96044 & 8024 & 8.4 & 59 & 38 & 193.8 \\
Cathlab & 9282 & 539 & 5.8 & 41.8 & 32 & 108.3 \\
Maaseik & 53902 & 1494 & 2.8 & 42 & 29 & 129 \\
Lanaken & 69395 & 371 & 0.5 & 15.7 & 10 & 48.5 \\
\end{longtable}

\newpage

\section{Are surgeries performed in their assigned operating
rooms?}\label{are-surgeries-performed-in-their-assigned-operating-rooms}

\subsection{How frequently are surgeries relocated to a different
room?}\label{how-frequently-are-surgeries-relocated-to-a-different-room}

Room reallocations are relatively rare at Lanaken, as shown in Table
\ref{tab:room_swap}. Out of 69395 recorded surgeries, around 7.5 percent
were performed in a different operating room than originally planned.
This low percentage indicates that room assignments are followed very
consistently, with only a small number of last-minute adjustments
required during daily operations.

Although infrequent, these relocations are operationally important
because they typically involve reactive scheduling decisions and
resource coordination. Even minor disruptions can have downstream
effects on room preparation, staff planning, and turnover time. Tracking
their frequency provides useful insight into the stability of the
planning process and helps identify whether reassignments occur
sporadically or tend to cluster in specific rooms or days.

\begin{longtable}[]{@{}
  >{\centering\arraybackslash}p{(\linewidth - 4\tabcolsep) * \real{0.1111}}
  >{\centering\arraybackslash}p{(\linewidth - 4\tabcolsep) * \real{0.2917}}
  >{\centering\arraybackslash}p{(\linewidth - 4\tabcolsep) * \real{0.1944}}@{}}
\caption{Overall frequency of surgeries relocated to a different room
(Lanaken). \label{tab:room_swap}}\tabularnewline
\toprule\noalign{}
\begin{minipage}[b]{\linewidth}\centering
Total
\end{minipage} & \begin{minipage}[b]{\linewidth}\centering
Swapped (absolute)
\end{minipage} & \begin{minipage}[b]{\linewidth}\centering
Swapped (\%)
\end{minipage} \\
\midrule\noalign{}
\endfirsthead
\toprule\noalign{}
\begin{minipage}[b]{\linewidth}\centering
Total
\end{minipage} & \begin{minipage}[b]{\linewidth}\centering
Swapped (absolute)
\end{minipage} & \begin{minipage}[b]{\linewidth}\centering
Swapped (\%)
\end{minipage} \\
\midrule\noalign{}
\endhead
\bottomrule\noalign{}
\endlastfoot
69395 & 5192 & 7.5 \\
\end{longtable}

The share of relocated surgeries remains low across all weekdays, as
shown in Figure \ref{fig:room_swap_weekday}, generally fluctuating
between 5.9 and 10.3 percent. This stability indicates that the room
assignment process is robust during regular working days, with only
minimal need for unplanned reallocations. The lowest rate is observed on
Wednesday, while Monday show higher values, of 10.3 percent.

\begin{figure}[H]

{\centering \includegraphics[width=0.7\linewidth,]{In-Depth_Analysis_Lanaken_files/figure-latex/room_swap_weekday-1} 

}

\caption{Share of relocated surgeries by weekday \label{fig:room_swap_weekday}}\label{fig:room_swap_weekday}
\end{figure}

The monthly percentage of surgeries that are relocated to a different
operating room fluctuates between 5 and 15.5 percent, as shown in Figure
\ref{fig:room_swap_month}. Despite short-term variation, no clear upward
or downward trend is visible over the observed period. Occasional peaks
above 10 percent indicate that some months experience more frequent
reallocations, but these appear as structural changes rather than
isolated events.

The relatively low and stable baseline confirms that room swaps are
infrequent and mainly driven by temporary scheduling pressures or
logistical constraints. The absence of a consistent trend over time
suggests that the overall stability of room assignments at Lanaken has
remained steady throughout the study period.

\begin{figure}[H]

{\centering \includegraphics[width=0.7\linewidth,]{In-Depth_Analysis_Lanaken_files/figure-latex/room_swap_month-1} 

}

\caption{Monthly percentage of surgeries with room swap \label{fig:room_swap_month}}\label{fig:room_swap_month}
\end{figure}

\subsection{Are specific rooms affected more than
others?}\label{are-specific-rooms-affected-more-than-others}

The analysis shows that only a limited number of operating rooms are
regularly involved in relocations, as summarized in Table
\ref{tab:room_swap_or}. A small subset of rooms accounts for nearly all
movements, while most others experience swaps only sporadically. The
most prominent interaction occurs between LP02 and LP01, which exchange
cases extensively: every surgery moved out of LP02 is reassigned to
LP01, and vice versa. This symmetry suggests that the two rooms function
as a tightly coupled pair with interchangeable capacity, absorbing
short-term adjustments when local scheduling pressure arises.

Within the main operating block, LO03 and LO04 display a similar
pattern, frequently receiving each other's relocated cases. LO03 also
sends a smaller share of cases to LO01 and LO02, indicating broader
compatibility with several rooms. LO05 participates in relocations as
well, though at a lower frequency, generally receiving overflow from
LO03 and LO04 rather than acting as a major source of swaps.

These consistent patterns show that reallocations follow clear
operational structures rather than occurring randomly. Rooms with
similar procedural profiles or shared staffing arrangements naturally
form relocation clusters. Movements therefore reflect adaptive use of
available capacity to maintain workflow continuity, rather than
instability in the planning process.

\begin{longtable}[]{@{}
  >{\centering\arraybackslash}p{(\linewidth - 10\tabcolsep) * \real{0.1165}}
  >{\centering\arraybackslash}p{(\linewidth - 10\tabcolsep) * \real{0.1262}}
  >{\centering\arraybackslash}p{(\linewidth - 10\tabcolsep) * \real{0.1942}}
  >{\centering\arraybackslash}p{(\linewidth - 10\tabcolsep) * \real{0.1942}}
  >{\centering\arraybackslash}p{(\linewidth - 10\tabcolsep) * \real{0.1845}}
  >{\centering\arraybackslash}p{(\linewidth - 10\tabcolsep) * \real{0.1845}}@{}}
\caption{Top 10 planned ORs with the highest number of room swaps,
showing the five most frequent new OR destinations (Lanaken).
\label{tab:room_swap_or}}\tabularnewline
\toprule\noalign{}
\begin{minipage}[b]{\linewidth}\centering
PlannedOR
\end{minipage} & \begin{minipage}[b]{\linewidth}\centering
TotalMoved
\end{minipage} & \begin{minipage}[b]{\linewidth}\centering
Top1
\end{minipage} & \begin{minipage}[b]{\linewidth}\centering
Top2
\end{minipage} & \begin{minipage}[b]{\linewidth}\centering
Top3
\end{minipage} & \begin{minipage}[b]{\linewidth}\centering
Top4
\end{minipage} \\
\midrule\noalign{}
\endfirsthead
\toprule\noalign{}
\begin{minipage}[b]{\linewidth}\centering
PlannedOR
\end{minipage} & \begin{minipage}[b]{\linewidth}\centering
TotalMoved
\end{minipage} & \begin{minipage}[b]{\linewidth}\centering
Top1
\end{minipage} & \begin{minipage}[b]{\linewidth}\centering
Top2
\end{minipage} & \begin{minipage}[b]{\linewidth}\centering
Top3
\end{minipage} & \begin{minipage}[b]{\linewidth}\centering
Top4
\end{minipage} \\
\midrule\noalign{}
\endhead
\bottomrule\noalign{}
\endlastfoot
LP02 & 2389 & LP01 (2389, 100\%) & NA & NA & NA \\
LP01 & 1140 & LP02 (1140, 100\%) & NA & NA & NA \\
LO03 & 588 & LO04 (248, 42.2\%) & LO05 (238, 40.5\%) & LO01 (93, 15.8\%)
& LO02 (9, 1.5\%) \\
LO05 & 431 & LO03 (242, 56.1\%) & LO04 (121, 28.1\%) & LO01 (65, 15.1\%)
& LO02 (3, 0.7\%) \\
LO04 & 358 & LO03 (216, 60.3\%) & LO05 (96, 26.8\%) & LO02 (37, 10.3\%)
& LO01 (9, 2.5\%) \\
LO01 & 184 & LO03 (94, 51.1\%) & LO05 (70, 38\%) & LO04 (13, 7.1\%) &
LO02 (7, 3.8\%) \\
LO02 & 100 & LO04 (87, 87\%) & LO03 (9, 9\%) & LO01 (2, 2\%) & LO05 (2,
2\%) \\
LOR0 & 2 & LO01 (1, 50\%) & LO05 (1, 50\%) & NA & NA \\
\end{longtable}

A small number of operating rooms receive the majority of relocated
surgeries, as shown in Table \ref{tab:room_swap_dest} and Figure
\ref{fig:room_swap_dest}. LP01 and LP02 stand out clearly: LP01 alone
accounts for 3.44 percent of all relocated cases, and LP02 adds another
1.64 percent. These two rooms function as the primary destinations when
schedule adjustments require shifting procedures, reflecting their
flexible capacity and broad procedural compatibility.

Several rooms in the main surgical block also appear regularly as
secondary destinations. LO03, LO04, and LO05 together receive a
nontrivial share of relocated cases, though far below the volumes seen
in LP01 and LP02. At the lower end of the distribution, rooms such as
LO01 and LO02 receive only a small number of relocated procedures.

Both the table and figure reveal a highly concentrated pattern: a small
subset of rooms absorbs most reallocations, while the remainder sees
minimal or occasional inflow. This concentration suggests that
relocation decisions are not random but guided by practical constraints
such as equipment availability, staffing configuration, and the capacity
of specific rooms to act as operational buffers.

\begin{longtable}[]{@{}
  >{\centering\arraybackslash}p{(\linewidth - 4\tabcolsep) * \real{0.1528}}
  >{\centering\arraybackslash}p{(\linewidth - 4\tabcolsep) * \real{0.2222}}
  >{\centering\arraybackslash}p{(\linewidth - 4\tabcolsep) * \real{0.1528}}@{}}
\caption{Top 15 actual operating rooms most frequently used as
destinations after relocation (Lanaken).
\label{tab:room_swap_dest}}\tabularnewline
\toprule\noalign{}
\begin{minipage}[b]{\linewidth}\centering
ActualOR
\end{minipage} & \begin{minipage}[b]{\linewidth}\centering
ReceivedCases
\end{minipage} & \begin{minipage}[b]{\linewidth}\centering
SharePct
\end{minipage} \\
\midrule\noalign{}
\endfirsthead
\toprule\noalign{}
\begin{minipage}[b]{\linewidth}\centering
ActualOR
\end{minipage} & \begin{minipage}[b]{\linewidth}\centering
ReceivedCases
\end{minipage} & \begin{minipage}[b]{\linewidth}\centering
SharePct
\end{minipage} \\
\midrule\noalign{}
\endhead
\bottomrule\noalign{}
\endlastfoot
LP01 & 2389 & 3.44 \\
LP02 & 1140 & 1.64 \\
LO03 & 561 & 0.81 \\
LO04 & 469 & 0.68 \\
LO05 & 407 & 0.59 \\
LO01 & 170 & 0.24 \\
LO02 & 56 & 0.08 \\
\end{longtable}

Procedures with the highest likelihood of being relocated mainly consist
of short or outpatient interventions, as shown in Table
\ref{tab:room_swap_procedure}. The top of the list is dominated by
minor, standardised procedures such as CARPAAL TUNNEL, Echo nervus
cutaneus femoralis lateralis PRF, and Echo perifere zenuw PRF, which
show relocation rates ranging from about 17 to 47 percent. These values
are several times higher than the overall relocation rate at Lanaken of
7.5 percent, indicating that a small subset of procedures accounts for a
disproportionate share of all room swaps.

These high-swap procedures share an important operational
characteristic: they are short, predictable, and compatible with
multiple procedure rooms. This makes them particularly suitable for
reassignment when minor scheduling conflicts or capacity pressures
arise. In contrast, longer or more specialised procedures rarely appear
among relocated cases, which reflects their greater dependence on
specific room configurations or equipment.

Overall, the table shows that relocations at Lanaken are concentrated in
a narrow group of routine, highly transferable procedures. This pattern
suggests that room swaps do not reflect instability in planning but
rather the deliberate use of flexible procedure types to absorb
day-to-day scheduling adjustments.

\begin{longtable}[]{@{}
  >{\centering\arraybackslash}p{(\linewidth - 4\tabcolsep) * \real{0.4583}}
  >{\centering\arraybackslash}p{(\linewidth - 4\tabcolsep) * \real{0.0972}}
  >{\centering\arraybackslash}p{(\linewidth - 4\tabcolsep) * \real{0.1389}}@{}}
\caption{Top 15 procedure types with the highest share of relocated
surgeries (Lanaken). \label{tab:room_swap_procedure}}\tabularnewline
\toprule\noalign{}
\begin{minipage}[b]{\linewidth}\centering
ProcedureName
\end{minipage} & \begin{minipage}[b]{\linewidth}\centering
n
\end{minipage} & \begin{minipage}[b]{\linewidth}\centering
SwapPct
\end{minipage} \\
\midrule\noalign{}
\endfirsthead
\toprule\noalign{}
\begin{minipage}[b]{\linewidth}\centering
ProcedureName
\end{minipage} & \begin{minipage}[b]{\linewidth}\centering
n
\end{minipage} & \begin{minipage}[b]{\linewidth}\centering
SwapPct
\end{minipage} \\
\midrule\noalign{}
\endhead
\bottomrule\noalign{}
\endlastfoot
CARPAAL TUNNEL & 1211 & 47.07 \\
Echo nervus cutaneus femoralis lateralis PRF & 52 & 42.31 \\
Echo perifere zenuw PRF & 73 & 38.36 \\
Echo bursa trochanterica & 170 & 29.41 \\
VARICES & 689 & 25.98 \\
Endscopische carpaal tunnel & 124 & 25.81 \\
NL RF facet lumbaal & 39 & 23.08 \\
Echo nervus genicularis RF (WOMAC) & 140 & 22.86 \\
VARICES LASER & 1351 & 22.65 \\
ZENUW CARPAAL TUNNEL & 495 & 21.21 \\
Ganglion sfenopalatinum : Radiofrequente R & 78 & 20.51 \\
HAND CYSTE & 62 & 19.35 \\
Echo buikwand/ACNES PRF & 32 & 18.75 \\
VINGER TRIGGERVINGER & 228 & 17.62 \\
Echo nervus ilio-inguinalis PRF & 35 & 17.14 \\
\end{longtable}

\subsection{What is the operational impact of these
reallocations?}\label{what-is-the-operational-impact-of-these-reallocations}

Relocated surgeries differ from non-relocated ones in several aspects of
timing and scheduling performance, as summarized in Table
\ref{tab:room_swap_impact}. For start times, relocated cases show a
higher average delay among late starts (39.4 minutes versus 31.7
minutes) and a larger lead among early starts (49.1 minutes versus 33.8
minutes). This suggests that these surgeries are more often scheduled
under atypical or time-constrained circumstances, which can result in
greater variability in when they actually begin. The median values,
however, remain closer together (26 versus 20 minutes for late starts,
36 versus 20 for early starts), indicating that while some relocated
surgeries start much later or earlier than expected, most are only
moderately affected.

For procedure duration, both relocated and non-relocated surgeries show
comparable deviations from their planned times. Average differences are
around 9 to 10 minutes for longer cases and 8 to 9 minutes for shorter
ones, suggesting that relocation itself does not substantially influence
how long a surgery takes once it begins.

Overtime remains limited in both groups, with median values at zero.
Mean overtime is identical between relocated and non-relocated cases
(0.1), which indicates that relocations rarely lead to extended working
hours.

Overall, the results indicate that relocations are associated with
greater variability in start timing but not with systematically longer
or shorter surgeries. Their operational impact therefore appears to be
primarily related to scheduling and coordination rather than to the
procedures themselves.

\begin{longtable}[]{@{}
  >{\centering\arraybackslash}p{(\linewidth - 6\tabcolsep) * \real{0.2222}}
  >{\centering\arraybackslash}p{(\linewidth - 6\tabcolsep) * \real{0.2639}}
  >{\centering\arraybackslash}p{(\linewidth - 6\tabcolsep) * \real{0.0972}}
  >{\centering\arraybackslash}p{(\linewidth - 6\tabcolsep) * \real{0.1250}}@{}}
\caption{Comparison of start and duration deviations (early/late,
shorter/longer) and overtime between relocated and non-relocated
surgeries (Lanaken). \label{tab:room_swap_impact}}\tabularnewline
\toprule\noalign{}
\begin{minipage}[b]{\linewidth}\centering
Type
\end{minipage} & \begin{minipage}[b]{\linewidth}\centering
Metric
\end{minipage} & \begin{minipage}[b]{\linewidth}\centering
Mean
\end{minipage} & \begin{minipage}[b]{\linewidth}\centering
Median
\end{minipage} \\
\midrule\noalign{}
\endfirsthead
\toprule\noalign{}
\begin{minipage}[b]{\linewidth}\centering
Type
\end{minipage} & \begin{minipage}[b]{\linewidth}\centering
Metric
\end{minipage} & \begin{minipage}[b]{\linewidth}\centering
Mean
\end{minipage} & \begin{minipage}[b]{\linewidth}\centering
Median
\end{minipage} \\
\midrule\noalign{}
\endhead
\bottomrule\noalign{}
\endlastfoot
Relocated & Late start & 39.4 & 26 \\
Not relocated & Late start & 31.7 & 21 \\
& ---- & & \\
Relocated & Early start & 49.1 & 36 \\
Not relocated & Early start & 33.8 & 20 \\
& ---- & & \\
Relocated & Longer duration & 9.9 & 6 \\
Not relocated & Longer duration & 8.9 & 5 \\
& ---- & & \\
Relocated & Shorter duration & 8.7 & 5 \\
Not relocated & Shorter duration & 7.8 & 5 \\
& ---- & & \\
Relocated & Overtime & 0.1 & 0 \\
Not relocated & Overtime & 0.1 & 0 \\
& ---- & & \\
\end{longtable}

Overtime is slightly more common among relocated surgeries than among
those performed in their originally assigned room, as shown in Table
\ref{tab:overtime_relocation}. About 0.5 percent of non-relocated
surgeries extend beyond 16:30, compared to 0.4 percent of relocated
ones. This small difference likely reflects that some relocated cases
are scheduled later in the day or in response to last-minute
adjustments, which can occasionally lead to minor overruns.

The similar percentages indicate that room reassignments are not a major
driver of evening work. Although relocations require coordination, they
do not substantially contribute to overtime or after-hours workload at
the campus level.

\begin{longtable}[]{@{}
  >{\centering\arraybackslash}p{(\linewidth - 4\tabcolsep) * \real{0.2222}}
  >{\centering\arraybackslash}p{(\linewidth - 4\tabcolsep) * \real{0.1806}}
  >{\centering\arraybackslash}p{(\linewidth - 4\tabcolsep) * \real{0.4306}}@{}}
\caption{Share of after-hours surgeries by relocation status (Lanaken).
\label{tab:overtime_relocation}}\tabularnewline
\toprule\noalign{}
\begin{minipage}[b]{\linewidth}\centering
SwapGroup
\end{minipage} & \begin{minipage}[b]{\linewidth}\centering
AfterHours
\end{minipage} & \begin{minipage}[b]{\linewidth}\centering
Percentage of surgeries with afterhours
\end{minipage} \\
\midrule\noalign{}
\endfirsthead
\toprule\noalign{}
\begin{minipage}[b]{\linewidth}\centering
SwapGroup
\end{minipage} & \begin{minipage}[b]{\linewidth}\centering
AfterHours
\end{minipage} & \begin{minipage}[b]{\linewidth}\centering
Percentage of surgeries with afterhours
\end{minipage} \\
\midrule\noalign{}
\endhead
\bottomrule\noalign{}
\endlastfoot
Not relocated & 348 & 0.5 \\
Relocated & 23 & 0.4 \\
\end{longtable}

The distribution of overtime duration for relocated surgeries, shown in
Figure \ref{fig:overtime_relocated}, is highly concentrated near 10 to
20, indicating that many of these cases finish slightly after regular
working hours or shortly afterward. A steep drop is visible beyond the
first 20 minutes of overtime, and only a few relocated surgeries extend
beyond 40 minutes The long but sparse tail confirms that extreme cases
are rare.

This pattern supports the earlier finding that relocations mostly
involve short procedures that can be completed quickly once a new room
becomes available. The limited overtime associated with these cases
suggests that room reassignments are generally managed efficiently and
do not significantly contribute to prolonged evening work. In
operational terms, relocated surgeries appear to represent flexible
workload adjustments rather than sources of additional delay.

\begin{figure}[H]

{\centering \includegraphics[width=0.7\linewidth,]{In-Depth_Analysis_Lanaken_files/figure-latex/overtime_relocated-1} 

}

\caption{Distribution of overtime duration for relocated surgeries \label{fig:overtime_relocated}}\label{fig:overtime_relocated}
\end{figure}

\newpage

\section{How do urgent cases disrupt elective
planning?}\label{how-do-urgent-cases-disrupt-elective-planning}

\subsection{When and where do urgent surgeries
occur?}\label{when-and-where-do-urgent-surgeries-occur}

The majority of surgeries at Lanaken are elective, representing about
99.6 percent of all recorded cases, as shown in Table
\ref{tab:urgency_dist}. Non-elective or urgent surgeries make up the
remaining 0.4 percent of the total. This indicates that most activity
follows a planned schedule, only a tiny portion of cases still arises
from unplanned or time-sensitive demands.

The relative size of the non-elective segment highlights its operational
significance. Even though urgent cases form a minority, their
unpredictable timing can affect the stability of daily planning and room
allocation. Understanding when and where these non-elective procedures
occur provides an important basis for evaluating how they interact with
elective activity and whether they contribute to schedule deviations or
after-hours work.

\begin{longtable}[]{@{}
  >{\centering\arraybackslash}p{(\linewidth - 4\tabcolsep) * \real{0.2222}}
  >{\centering\arraybackslash}p{(\linewidth - 4\tabcolsep) * \real{0.1111}}
  >{\centering\arraybackslash}p{(\linewidth - 4\tabcolsep) * \real{0.1528}}@{}}
\caption{Distribution of surgeries by urgency classification (Lanaken).
\label{tab:urgency_dist}}\tabularnewline
\toprule\noalign{}
\begin{minipage}[b]{\linewidth}\centering
UrgencyType
\end{minipage} & \begin{minipage}[b]{\linewidth}\centering
Total
\end{minipage} & \begin{minipage}[b]{\linewidth}\centering
SharePct
\end{minipage} \\
\midrule\noalign{}
\endfirsthead
\toprule\noalign{}
\begin{minipage}[b]{\linewidth}\centering
UrgencyType
\end{minipage} & \begin{minipage}[b]{\linewidth}\centering
Total
\end{minipage} & \begin{minipage}[b]{\linewidth}\centering
SharePct
\end{minipage} \\
\midrule\noalign{}
\endhead
\bottomrule\noalign{}
\endlastfoot
electief & 69109 & 99.6 \\
niet-electief & 286 & 0.4 \\
\end{longtable}

The start times of non-elective surgeries, shown in Figure
\ref{fig:urgency_timing}, display a clear concentration during daytime
hours, with a pronounced peak between 13:00 and 15:00. This pattern
indicates that the majority of urgent procedures are performed within or
near regular operating hours, often overlapping with scheduled elective
activity.

\begin{figure}[H]

{\centering \includegraphics[width=0.7\linewidth,]{In-Depth_Analysis_Lanaken_files/figure-latex/urgency_timing-1} 

}

\caption{Daily distribution of non-elective surgery start times \label{fig:urgency_timing}}\label{fig:urgency_timing}
\end{figure}

Across all weekdays, non-elective surgeries follow a consistent hourly
pattern, as illustrated in Figure \ref{fig:urgency_weekday}. Most
activity occurs between late morning and early evening, with peaks
around 13:00 to 15:00 on nearly every day. This confirms that urgent
procedures are frequently performed during regular daytime hours. The
overlap with elective schedules suggests that urgent cases often need to
be accommodated within existing capacity rather than outside normal
shifts.

\begin{figure}[H]

{\centering \includegraphics[width=0.7\linewidth,]{In-Depth_Analysis_Lanaken_files/figure-latex/urgency_weekday-1} 

}

\caption{Hourly distribution of non-elective surgery start times by weekday \label{fig:urgency_weekday}}\label{fig:urgency_weekday}
\end{figure}

Elective surgeries dominate the activity, as shown in Table
\ref{tab:urgency_wd}, accounting for roughly 99 percent of all cases
from Monday to Friday. The share of non-elective procedures remains
stable during the workweek, between 0.2 and 0.8 percent, indicating a
consistent and predictable pattern of urgent demand. This stability
suggests that urgent cases are routinely managed alongside scheduled
operations without major fluctuations across days.

\begin{longtable}[]{@{}
  >{\centering\arraybackslash}p{(\linewidth - 8\tabcolsep) * \real{0.1282}}
  >{\centering\arraybackslash}p{(\linewidth - 8\tabcolsep) * \real{0.1923}}
  >{\centering\arraybackslash}p{(\linewidth - 8\tabcolsep) * \real{0.2436}}
  >{\centering\arraybackslash}p{(\linewidth - 8\tabcolsep) * \real{0.1923}}
  >{\centering\arraybackslash}p{(\linewidth - 8\tabcolsep) * \real{0.2436}}@{}}
\caption{Weekday distribution of elective vs.~non-elective surgeries
(Lanaken). \label{tab:urgency_wd}}\tabularnewline
\toprule\noalign{}
\begin{minipage}[b]{\linewidth}\centering
Weekday
\end{minipage} & \begin{minipage}[b]{\linewidth}\centering
Elective (n)
\end{minipage} & \begin{minipage}[b]{\linewidth}\centering
Non-elective (n)
\end{minipage} & \begin{minipage}[b]{\linewidth}\centering
Elective (\%)
\end{minipage} & \begin{minipage}[b]{\linewidth}\centering
Non-elective (\%)
\end{minipage} \\
\midrule\noalign{}
\endfirsthead
\toprule\noalign{}
\begin{minipage}[b]{\linewidth}\centering
Weekday
\end{minipage} & \begin{minipage}[b]{\linewidth}\centering
Elective (n)
\end{minipage} & \begin{minipage}[b]{\linewidth}\centering
Non-elective (n)
\end{minipage} & \begin{minipage}[b]{\linewidth}\centering
Elective (\%)
\end{minipage} & \begin{minipage}[b]{\linewidth}\centering
Non-elective (\%)
\end{minipage} \\
\midrule\noalign{}
\endhead
\bottomrule\noalign{}
\endlastfoot
Ma & 14635 & 75 & 99.5 & 0.5 \\
Di & 14878 & 24 & 99.8 & 0.2 \\
Wo & 13683 & 59 & 99.6 & 0.4 \\
Do & 13209 & 24 & 99.8 & 0.2 \\
Vr & 12704 & 104 & 99.2 & 0.8 \\
\end{longtable}

\subsection{Do urgent cases drive
overtime?}\label{do-urgent-cases-drive-overtime}

Urgent cases contribute to overtime work but are not its primary driver,
as shown in Table \ref{tab:afterhours_urgency}. While the share of
non-elective surgeries extending beyond the shift is higher (3.8
percent) than for elective ones (0.5 percent), the overall difference
remains moderate. This indicates that both types of activity continue
after the shift, although elective cases represent the majority in
absolute numbers due to their much larger total volume.

Average overtime per case is slightly higher for non-elective surgeries,
at 0.4 minutes compared with 0.1 minutes for elective ones. The 95th
percentile follows a different pattern, reaching 0 minutes for both
urgent cases and for elective ones. This suggests that urgent
procedures, much less frequent, barely require longer extensions when
they do occur.

Overall, the results imply that unplanned demand does contribute to
extended operating hours, but most work still stems from elective cases
that overrun their planned schedules rather than from late-arriving
emergencies.

\begin{longtable}[]{@{}
  >{\centering\arraybackslash}p{(\linewidth - 10\tabcolsep) * \real{0.1376}}
  >{\centering\arraybackslash}p{(\linewidth - 10\tabcolsep) * \real{0.1101}}
  >{\centering\arraybackslash}p{(\linewidth - 10\tabcolsep) * \real{0.1651}}
  >{\centering\arraybackslash}p{(\linewidth - 10\tabcolsep) * \real{0.1651}}
  >{\centering\arraybackslash}p{(\linewidth - 10\tabcolsep) * \real{0.2018}}
  >{\centering\arraybackslash}p{(\linewidth - 10\tabcolsep) * \real{0.2202}}@{}}
\caption{After-hours frequency and overtime duration by urgency type
(Lanaken). \label{tab:afterhours_urgency}}\tabularnewline
\toprule\noalign{}
\begin{minipage}[b]{\linewidth}\centering
Urgency type
\end{minipage} & \begin{minipage}[b]{\linewidth}\centering
Total (n)
\end{minipage} & \begin{minipage}[b]{\linewidth}\centering
After-hours (n)
\end{minipage} & \begin{minipage}[b]{\linewidth}\centering
After-hours (\%)
\end{minipage} & \begin{minipage}[b]{\linewidth}\centering
Mean overtime (min)
\end{minipage} & \begin{minipage}[b]{\linewidth}\centering
95th percentile (min)
\end{minipage} \\
\midrule\noalign{}
\endfirsthead
\toprule\noalign{}
\begin{minipage}[b]{\linewidth}\centering
Urgency type
\end{minipage} & \begin{minipage}[b]{\linewidth}\centering
Total (n)
\end{minipage} & \begin{minipage}[b]{\linewidth}\centering
After-hours (n)
\end{minipage} & \begin{minipage}[b]{\linewidth}\centering
After-hours (\%)
\end{minipage} & \begin{minipage}[b]{\linewidth}\centering
Mean overtime (min)
\end{minipage} & \begin{minipage}[b]{\linewidth}\centering
95th percentile (min)
\end{minipage} \\
\midrule\noalign{}
\endhead
\bottomrule\noalign{}
\endlastfoot
Elective & 69109 & 360 & 0.5 & 0.1 & 0 \\
NonElective & 286 & 11 & 3.8 & 0.4 & 0 \\
\end{longtable}

\subsection{Do urgent cases interrupt elective
planning?}\label{do-urgent-cases-interrupt-elective-planning}

The results in Table \ref{tab:overlap_swap} show that the frequency of
room swaps among elective surgeries is virtually identical regardless of
whether an overlap with a non-elective case occurred. For surgeries not
affected by an overlap, 7.5 percent involved a room swap, compared to
the 3.5 percent among overlapped cases.

This consistency indicates that while urgent procedures occasionally
coincide with scheduled elective activity, they do not systematically
trigger room reallocations. Instead, such overlaps appear to be managed
through local coordination or minor timing adjustments rather than by
physically moving elective cases to different operating rooms. Overall,
urgent interventions create some scheduling pressure but do not
significantly disrupt room allocation patterns within the elective
program.

\begin{longtable}[]{@{}
  >{\centering\arraybackslash}p{(\linewidth - 6\tabcolsep) * \real{0.1806}}
  >{\centering\arraybackslash}p{(\linewidth - 6\tabcolsep) * \real{0.1111}}
  >{\centering\arraybackslash}p{(\linewidth - 6\tabcolsep) * \real{0.1667}}
  >{\centering\arraybackslash}p{(\linewidth - 6\tabcolsep) * \real{0.2083}}@{}}
\caption{Room swap frequency among elective cases: overlapped vs.~not
overlapped. \label{tab:overlap_swap}}\tabularnewline
\toprule\noalign{}
\begin{minipage}[b]{\linewidth}\centering
Overlapped
\end{minipage} & \begin{minipage}[b]{\linewidth}\centering
N
\end{minipage} & \begin{minipage}[b]{\linewidth}\centering
RoomSwaps
\end{minipage} & \begin{minipage}[b]{\linewidth}\centering
ShareSwapped
\end{minipage} \\
\midrule\noalign{}
\endfirsthead
\toprule\noalign{}
\begin{minipage}[b]{\linewidth}\centering
Overlapped
\end{minipage} & \begin{minipage}[b]{\linewidth}\centering
N
\end{minipage} & \begin{minipage}[b]{\linewidth}\centering
RoomSwaps
\end{minipage} & \begin{minipage}[b]{\linewidth}\centering
ShareSwapped
\end{minipage} \\
\midrule\noalign{}
\endhead
\bottomrule\noalign{}
\endlastfoot
No & 68652 & 5169 & 7.5 \\
Yes & 455 & 16 & 3.5 \\
\end{longtable}

As shown in Table \ref{tab:overlap_days}, overlaps between urgent and
elective surgeries occurred on 143 out of 853 observed days,
corresponding to 16.8 percent of all days. This moderate frequency
indicates that schedule interference is not an exceptional event but a
regular characteristic of daily activity. While individual overlaps may
vary in magnitude and operational impact, their near-daily occurrence
highlights the continuous need for flexibility in resource allocation
and real-time coordination between planned and urgent care.

\begin{longtable}[]{@{}
  >{\centering\arraybackslash}p{(\linewidth - 4\tabcolsep) * \real{0.1944}}
  >{\centering\arraybackslash}p{(\linewidth - 4\tabcolsep) * \real{0.1667}}
  >{\centering\arraybackslash}p{(\linewidth - 4\tabcolsep) * \real{0.1667}}@{}}
\caption{Days with at least one urgent--elective overlap in the same OR.
\label{tab:overlap_days}}\tabularnewline
\toprule\noalign{}
\begin{minipage}[b]{\linewidth}\centering
OverlapDays
\end{minipage} & \begin{minipage}[b]{\linewidth}\centering
TotalDays
\end{minipage} & \begin{minipage}[b]{\linewidth}\centering
SharePct
\end{minipage} \\
\midrule\noalign{}
\endfirsthead
\toprule\noalign{}
\begin{minipage}[b]{\linewidth}\centering
OverlapDays
\end{minipage} & \begin{minipage}[b]{\linewidth}\centering
TotalDays
\end{minipage} & \begin{minipage}[b]{\linewidth}\centering
SharePct
\end{minipage} \\
\midrule\noalign{}
\endhead
\bottomrule\noalign{}
\endlastfoot
143 & 853 & 16.8 \\
\end{longtable}

The proportion of elective surgeries affected by urgent overlaps
fluctuates modestly over time, as shown in Figure
\ref{fig:overlap_trend}. Monthly values generally range between 0.4 and
0.8 percent. A gradual upward trend is visible throughout the end of
2023 and 2024. Despite these month-to-month variations, the overall
level remains relatively stable, suggesting that the frequency of urgent
cases interfering with scheduled activity is a persistent feature of
daily operations rather than a short-term anomaly.

\begin{figure}[H]

{\centering \includegraphics[width=0.7\linewidth,]{In-Depth_Analysis_Lanaken_files/figure-latex/overlap_trend-1} 

}

\caption{Monthly share of elective surgeries affected by urgent overlaps \label{fig:overlap_trend}}\label{fig:overlap_trend}
\end{figure}

As shown in Figure \ref{fig:overlap_hourly}, overlaps between urgent and
elective surgeries occur during daytime hours. The frequency increases
sharply after 12:00, peaks between 13:00 and 15:00, and then declines
sharply toward 16:00. The clustering of overlaps in the middle of the
day suggests that urgent admissions frequently coincide with the busiest
elective periods, intensifying competition for operating room capacity
during standard working hours.

\begin{figure}[H]

{\centering \includegraphics[width=0.7\linewidth,]{In-Depth_Analysis_Lanaken_files/figure-latex/overlap_hourly-1} 

}

\caption{Hourly frequency of urgent–elective overlaps \label{fig:overlap_hourly}}\label{fig:overlap_hourly}
\end{figure}

Figure \ref{fig:overlap_delay} shows clear differences in average start
delays between elective surgeries with and without an urgent overlap.
During much of 2022, the overlapped cases display strongly negative mean
delays, in several months exceeding minus 50 minutes, which indicates
that these surgeries tended to start significantly earlier than planned.
In comparison, non-overlapped cases remained close to the planned
schedule, with mean delays typically around zero to ten minutes.

From 2023 onward the gap between both groups narrows, yet overlapped
cases continue to fluctuate more widely, with repeated downward spikes
and intermittent positive peaks. This increased volatility suggests that
the timing of surgeries affected by urgent activity is more sensitive to
local workflow pressures, leading them at times to start markedly
earlier or slightly later than planned.

The recurring peaks and troughs in the overlapped series indicate that
overlaps are associated with periodic scheduling tensions rather than
isolated irregularities. Overall, the figure shows that urgent cases
introduce measurable variability into elective scheduling, not through
systematic delays but through greater unpredictability in start times,
especially during periods of high operating room utilisation.

\begin{figure}[H]

{\centering \includegraphics[width=0.7\linewidth,]{In-Depth_Analysis_Lanaken_files/figure-latex/overlap_delay-1} 

}

\caption{Average start delay of elective surgeries by overlap status \label{fig:overlap_delay}}\label{fig:overlap_delay}
\end{figure}

As shown in Table \ref{tab:or_overlap}, the frequency and impact of
urgent--elective overlaps vary substantially across operating rooms. The
room with the highest number of overlap events is LO02, with 259
overlapped cases, representing 2.6 percent of its elective activity.
This confirms that LO02 is one of the primary locations where urgent
demand intersects with scheduled work. Other rooms with notable counts
include LO04 (108 overlaps), LP01 (33), and LO03 (19), although their
relative shares remain below 2 percent.

Median start-delay patterns show that overlaps do not affect all rooms
in the same way. In LO04 and LO03, the median delay under overlap
becomes negative (--16.5 and --7 minutes respectively), indicating that
overlapped elective cases tend to start earlier than planned. LO01 shows
a similar pattern with a median of --3 minutes. In contrast, rooms such
as LP01 and LP02 exhibit positive median delays under overlap (12 and 9
minutes), suggesting that overlaps in these rooms are more likely to
push elective cases later than scheduled.

A few rooms such as LO05 display extreme values (--73 minutes), but
because these rooms have very small numbers of overlaps, these medians
should be interpreted with caution.

Overall, the results show that the operational effect of
urgent--elective overlap is concentrated in a small number of
high-volume rooms. Even within these rooms, the direction of the impact
is not uniform: in some settings overlaps delay elective cases, while in
others they coincide with earlier-than-planned starts, likely reflecting
ad hoc rescheduling, case reordering, or opportunistic use of available
capacity.

\begin{longtable}[]{@{}
  >{\centering\arraybackslash}p{(\linewidth - 10\tabcolsep) * \real{0.1165}}
  >{\centering\arraybackslash}p{(\linewidth - 10\tabcolsep) * \real{0.0777}}
  >{\centering\arraybackslash}p{(\linewidth - 10\tabcolsep) * \real{0.1262}}
  >{\centering\arraybackslash}p{(\linewidth - 10\tabcolsep) * \real{0.1262}}
  >{\centering\arraybackslash}p{(\linewidth - 10\tabcolsep) * \real{0.2913}}
  >{\centering\arraybackslash}p{(\linewidth - 10\tabcolsep) * \real{0.2621}}@{}}
\caption{Operating rooms with urgent--elective overlap.
\label{tab:or_overlap}}\tabularnewline
\toprule\noalign{}
\begin{minipage}[b]{\linewidth}\centering
PlannedOR
\end{minipage} & \begin{minipage}[b]{\linewidth}\centering
Count
\end{minipage} & \begin{minipage}[b]{\linewidth}\centering
n\_elective
\end{minipage} & \begin{minipage}[b]{\linewidth}\centering
Percentage
\end{minipage} & \begin{minipage}[b]{\linewidth}\centering
Med StartDelay (No Overlap)
\end{minipage} & \begin{minipage}[b]{\linewidth}\centering
Med StartDelay (Overlap)
\end{minipage} \\
\midrule\noalign{}
\endfirsthead
\toprule\noalign{}
\begin{minipage}[b]{\linewidth}\centering
PlannedOR
\end{minipage} & \begin{minipage}[b]{\linewidth}\centering
Count
\end{minipage} & \begin{minipage}[b]{\linewidth}\centering
n\_elective
\end{minipage} & \begin{minipage}[b]{\linewidth}\centering
Percentage
\end{minipage} & \begin{minipage}[b]{\linewidth}\centering
Med StartDelay (No Overlap)
\end{minipage} & \begin{minipage}[b]{\linewidth}\centering
Med StartDelay (Overlap)
\end{minipage} \\
\midrule\noalign{}
\endhead
\bottomrule\noalign{}
\endlastfoot
LO02 & 259 & 9871 & 2.6 & 6 & 0 \\
LO04 & 108 & 5806 & 1.9 & 13 & -16.5 \\
LP01 & 33 & 18936 & 0.2 & 10 & 12 \\
LO03 & 19 & 4206 & 0.5 & 5 & -7 \\
LO01 & 17 & 5200 & 0.3 & 0 & -3 \\
LP02 & 16 & 19863 & 0.1 & -2 & 9 \\
LO05 & 3 & 5223 & 0.1 & 5 & -73 \\
\end{longtable}

\newpage

\section{How stable is the surgical schedule in
execution?}\label{how-stable-is-the-surgical-schedule-in-execution}

Beyond accurate planning, the stability of the realised schedule is a
key indicator of operational performance in the operating theatre. Even
when planned start times, durations, and room assignments are correct,
day-to-day execution may deviate because of unforeseen delays, emergency
insertions, or coordination issues.

This section evaluates how consistently the planned schedule is
maintained throughout the day by focusing on two complementary aspects:

\begin{itemize}
\tightlist
\item
  \emph{Shift changes}, where surgeries are executed in a different
  shift than originally planned, indicating how often the daily
  programme must be adjusted in real time.
\item
  \emph{Idle or gap times}, representing the pauses between consecutive
  surgeries within the same room, which reflect how efficiently the
  available capacity is utilised.
\end{itemize}

Together, these indicators capture the operational resilience of the OR
complex and the ability to absorb disruptions and maintain smooth
throughput despite variable circumstances. High rates of shift changes
or long idle periods point to latent inefficiencies in daily planning
and coordination, whereas stability across these metrics signals a
well-synchronised workflow.

\subsection{How often are surgeries moved to another
shift?}\label{how-often-are-surgeries-moved-to-another-shift}

Changes between planned and actual shifts occur when surgeries start
significantly later than scheduled or are postponed to a later block of
the day. The actual shift is determined based on the observed start time
of each surgery relative to the official shift boundaries used at the
operating department (day: 08:00--16:30, evening: 16:30--22:00, night:
22:00--08:00).

At Campus Lanaken, 578 surgeries, or 0.8 percent of all recorded cases,
were performed in a different shift than originally planned, as shown in
Table \ref{tab:shift_mismatch}. For these surgeries, the average start
delay was minus 143.1 minutes, indicating that such deviations typically
involve a substantial advance within the daily schedule. Interestingly,
the average deviation in procedure duration was minus 11.7 minutes,
meaning that surgeries executed in another shift tend to be somewhat
shorter than their planned estimates once they begin.

The mean overtime for these cases was 1 minute, showing that some shift
changes result from late starts or, in other cases, procedures that
exceed their planned duration, rather than extensive spillovers deep
into the next shift. Although these cases represent only a small
fraction of total surgical activity, their magnitude and timing can
still disrupt staffing and sequencing for subsequent sessions.

\begin{longtable}[]{@{}
  >{\centering\arraybackslash}p{(\linewidth - 10\tabcolsep) * \real{0.0690}}
  >{\centering\arraybackslash}p{(\linewidth - 10\tabcolsep) * \real{0.1379}}
  >{\centering\arraybackslash}p{(\linewidth - 10\tabcolsep) * \real{0.1552}}
  >{\centering\arraybackslash}p{(\linewidth - 10\tabcolsep) * \real{0.2155}}
  >{\centering\arraybackslash}p{(\linewidth - 10\tabcolsep) * \real{0.2328}}
  >{\centering\arraybackslash}p{(\linewidth - 10\tabcolsep) * \real{0.1897}}@{}}
\caption{Cases performed in a different shift than planned (Lanaken).
\label{tab:shift_mismatch}}\tabularnewline
\toprule\noalign{}
\begin{minipage}[b]{\linewidth}\centering
n
\end{minipage} & \begin{minipage}[b]{\linewidth}\centering
MismatchCount
\end{minipage} & \begin{minipage}[b]{\linewidth}\centering
Share moved (\%)
\end{minipage} & \begin{minipage}[b]{\linewidth}\centering
Mean Start Delay (min)
\end{minipage} & \begin{minipage}[b]{\linewidth}\centering
Mean Duration Diff (min)
\end{minipage} & \begin{minipage}[b]{\linewidth}\centering
Mean Overtime (min)
\end{minipage} \\
\midrule\noalign{}
\endfirsthead
\toprule\noalign{}
\begin{minipage}[b]{\linewidth}\centering
n
\end{minipage} & \begin{minipage}[b]{\linewidth}\centering
MismatchCount
\end{minipage} & \begin{minipage}[b]{\linewidth}\centering
Share moved (\%)
\end{minipage} & \begin{minipage}[b]{\linewidth}\centering
Mean Start Delay (min)
\end{minipage} & \begin{minipage}[b]{\linewidth}\centering
Mean Duration Diff (min)
\end{minipage} & \begin{minipage}[b]{\linewidth}\centering
Mean Overtime (min)
\end{minipage} \\
\midrule\noalign{}
\endhead
\bottomrule\noalign{}
\endlastfoot
69395 & 578 & 0.8 & -143.1 & -11.7 & 1 \\
\end{longtable}

As shown in Figure \ref{fig:shift_moved_monthly}, the monthly share of
surgeries performed in a different shift oscillates visibly but without
any structural upward or downward trend. In early 2022, values generally
range between about 0.8 and 1.3 percent, with a small number of peaks
reaching roughly 1.5 percent. Throughout the remainder of the period,
the pattern continues to fluctuate between approximately 0.3 and 1.3
percent, reflecting normal month-to-month variation rather than
systematic change. Toward 2025, the series appears to settle into a
narrower band around 0.7 to 0.9 percent.

Overall, the figure indicates that shift changes occur regularly but at
a consistently low rate, and their frequency does not show signs of
long-term increase or deterioration over time.

\begin{figure}[H]

{\centering \includegraphics[width=0.7\linewidth,]{In-Depth_Analysis_Lanaken_files/figure-latex/shift_moved_monthly-1} 

}

\caption{Monthly percentage of surgeries moved to another shift \label{fig:shift_moved_monthly}}\label{fig:shift_moved_monthly}
\end{figure}

Overall, the data suggest a modest improvement in schedule adherence
after the fluctuations seen in early 2022, followed by a period of
relative stability. The recurring month-to-month peaks and dips appear
to reflect normal variation in workload, staffing, and case mix rather
than underlying structural scheduling issues. The absence of a
persistent upward trend indicates that shift changes remain a controlled
and predictable feature of daily operations.

\subsection{How much idle time occurs between consecutive
surgeries?}\label{how-much-idle-time-occurs-between-consecutive-surgeries}

Idle or gap time represents the short pause between the end of one
surgery and the start of the next within the same operating room during
the regular day shift (08:00--16:30). It captures how efficiently rooms
are prepared for the next case and how well surgical schedules are
synchronised in daily practice. For each room-day combination, the gap
time measures either the delay between shift start and the first case or
the interval between consecutive cases, excluding breaks longer than 60
minutes that reflect planned downtime.

At Campus Lanaken, idle time between consecutive surgeries is extremely
low. As shown in Table \ref{tab:idle_summary}, the average gap is 3.2
minutes, but the median is 0 minutes, meaning that in at least half of
all transitions the next case starts immediately after the previous one
ends. The interquartile range of 3 minutes confirms that most turnover
intervals fall within a very narrow band, reflecting tightly sequenced
workflows and minimal unused room time.

Longer gaps occur only infrequently. The 95th percentile reaches 15
minutes, and only the extreme tail approaches the cutoff of 60 minutes,
which is used to exclude planned breaks rather than true idle time.
Overall, the distribution indicates a highly efficient turnover process,
with very limited room inactivity between cases.

\begin{longtable}[]{@{}
  >{\centering\arraybackslash}p{(\linewidth - 16\tabcolsep) * \real{0.0972}}
  >{\centering\arraybackslash}p{(\linewidth - 16\tabcolsep) * \real{0.1250}}
  >{\centering\arraybackslash}p{(\linewidth - 16\tabcolsep) * \real{0.0833}}
  >{\centering\arraybackslash}p{(\linewidth - 16\tabcolsep) * \real{0.0833}}
  >{\centering\arraybackslash}p{(\linewidth - 16\tabcolsep) * \real{0.0833}}
  >{\centering\arraybackslash}p{(\linewidth - 16\tabcolsep) * \real{0.0833}}
  >{\centering\arraybackslash}p{(\linewidth - 16\tabcolsep) * \real{0.0833}}
  >{\centering\arraybackslash}p{(\linewidth - 16\tabcolsep) * \real{0.0833}}
  >{\centering\arraybackslash}p{(\linewidth - 16\tabcolsep) * \real{0.1111}}@{}}
\caption{Summary statistics of idle time between consecutive surgeries
(minutes). \label{tab:idle_summary}}\tabularnewline
\toprule\noalign{}
\begin{minipage}[b]{\linewidth}\centering
Mean
\end{minipage} & \begin{minipage}[b]{\linewidth}\centering
Median
\end{minipage} & \begin{minipage}[b]{\linewidth}\centering
SD
\end{minipage} & \begin{minipage}[b]{\linewidth}\centering
IQR
\end{minipage} & \begin{minipage}[b]{\linewidth}\centering
Min
\end{minipage} & \begin{minipage}[b]{\linewidth}\centering
Max
\end{minipage} & \begin{minipage}[b]{\linewidth}\centering
P95
\end{minipage} & \begin{minipage}[b]{\linewidth}\centering
P99
\end{minipage} & \begin{minipage}[b]{\linewidth}\centering
P99.9
\end{minipage} \\
\midrule\noalign{}
\endfirsthead
\toprule\noalign{}
\begin{minipage}[b]{\linewidth}\centering
Mean
\end{minipage} & \begin{minipage}[b]{\linewidth}\centering
Median
\end{minipage} & \begin{minipage}[b]{\linewidth}\centering
SD
\end{minipage} & \begin{minipage}[b]{\linewidth}\centering
IQR
\end{minipage} & \begin{minipage}[b]{\linewidth}\centering
Min
\end{minipage} & \begin{minipage}[b]{\linewidth}\centering
Max
\end{minipage} & \begin{minipage}[b]{\linewidth}\centering
P95
\end{minipage} & \begin{minipage}[b]{\linewidth}\centering
P99
\end{minipage} & \begin{minipage}[b]{\linewidth}\centering
P99.9
\end{minipage} \\
\midrule\noalign{}
\endhead
\bottomrule\noalign{}
\endlastfoot
3.2 & 0 & 7.8 & 3 & 0 & 60 & 15 & 46 & 59 \\
\end{longtable}

As shown in Figure \ref{fig:idle_distribution}, the histogram confirms
that idle intervals at Campus Lanaken are overwhelmingly short. The
first bar dominates the distribution: the vast majority of all
transitions fall below 5 minutes, indicating that turnover often occurs
almost immediately after the previous case finishes. The next few bins
show a rapid decline, with only a small fraction of intervals extending
into the 5--10 or 10--15 minute range.

Beyond 15 minutes, the frequency drops sharply, and the right-hand tail
becomes extremely thin. Only isolated transitions reach higher values,
and gaps longer than one hour are excluded by design.

Overall, the distribution illustrates a highly efficient turnover
process. Idle periods between surgeries are consistently brief, and
longer gaps, while present, are rare and do not represent a systematic
pattern of underutilisation.

\begin{figure}[H]

{\centering \includegraphics[width=0.7\linewidth,]{In-Depth_Analysis_Lanaken_files/figure-latex/idle_distribution-1} 

}

\caption{Distribution of idle time between surgeries (grouped per 5 min) \label{fig:idle_distribution}}\label{fig:idle_distribution}
\end{figure}

\subsection{How does idle time vary across operating
rooms?}\label{how-does-idle-time-vary-across-operating-rooms}

As shown in Figure \ref{fig:idle_or}, idle times vary noticeably across
operating rooms at Campus Lanaken. In every room, the median idle time
is considerably lower than the mean, which reflects the strongly
right-skewed nature of turnover intervals: most transitions occur with
little to no idle time, while a smaller number of longer gaps pull the
average upward.

The room-to-room differences are substantial. LO03 stands out with the
highest idle times, showing both a large mean and a clearly elevated
median, indicating that this room experiences longer and more frequent
pauses between cases. LO05 and LO01 also show higher mean idle times
compared with the rest of the OR complex, though their medians remain
relatively low, suggesting occasional extended gaps rather than
systematically slower turnover.

In contrast, rooms such as LP02, LP01, and especially LO02 display very
short idle times, with medians at or near zero and relatively low means.
These rooms appear to achieve consistently rapid turnover, likely
reflecting high-throughput workflows or more predictable case sequences.

Overall, the figure shows that while turnover efficiency is generally
strong across Lanaken, a few rooms exhibit noticeably longer idle
intervals. These differences point to room-specific scheduling patterns
or case mixes rather than systemic inefficiency.

\begin{figure}[H]

{\centering \includegraphics[width=0.7\linewidth,]{In-Depth_Analysis_Lanaken_files/figure-latex/idle_or-1} 

}

\caption{Mean and median idle time per operating room \label{fig:idle_or}}\label{fig:idle_or}
\end{figure}

As shown in Table \ref{tab:idle_variability}, the standard deviation
(SD) and coefficient of variation (CV) summarise how much idle times
fluctuate around their typical value for each operating room. The SD
ranges from about 5 to 11 minutes across rooms, confirming that most
turnover intervals remain short even when variability increases. The CV
values, however, are very high---between roughly 130 and 344
percent---because mean and median idle times in Lanaken are extremely
small, causing even modest fluctuations to translate into large relative
variation.

Rooms such as LO03, LO05, and LO04 stand out with the highest SDs,
indicating broader spreads in turnover times and less predictable
intervals between cases. These rooms experience longer gaps more
frequently, consistent with what is seen in the mean--median patterns in
Figure \ref{fig:idle_or}.

Toward the bottom of the table, rooms such as LP01, LO02, and LP02
display very high CVs despite relatively low SDs. This arises from the
ratio-based nature of the CV: when typical idle times are extremely
short---as they are in these rooms---even small fluctuations produce a
disproportionately large coefficient of variation.

Overall, variability in idle time differs meaningfully between rooms,
but the underlying distributions remain dominated by very short turnover
intervals. Only a small subset of rooms shows broader spreads that may
reflect more diverse schedules, irregular case sequencing, or
room-specific workflow constraints.

\begin{longtable}[]{@{}
  >{\centering\arraybackslash}p{(\linewidth - 6\tabcolsep) * \real{0.1528}}
  >{\centering\arraybackslash}p{(\linewidth - 6\tabcolsep) * \real{0.1111}}
  >{\centering\arraybackslash}p{(\linewidth - 6\tabcolsep) * \real{0.2917}}
  >{\centering\arraybackslash}p{(\linewidth - 6\tabcolsep) * \real{0.3750}}@{}}
\caption{Standard deviation and coefficient of variation of idle time
per operating room (minutes), sorted by median idle time.
\label{tab:idle_variability}}\tabularnewline
\toprule\noalign{}
\begin{minipage}[b]{\linewidth}\centering
ActualOR
\end{minipage} & \begin{minipage}[b]{\linewidth}\centering
Cases
\end{minipage} & \begin{minipage}[b]{\linewidth}\centering
Standard deviation
\end{minipage} & \begin{minipage}[b]{\linewidth}\centering
Coefficient of variation
\end{minipage} \\
\midrule\noalign{}
\endfirsthead
\toprule\noalign{}
\begin{minipage}[b]{\linewidth}\centering
ActualOR
\end{minipage} & \begin{minipage}[b]{\linewidth}\centering
Cases
\end{minipage} & \begin{minipage}[b]{\linewidth}\centering
Standard deviation
\end{minipage} & \begin{minipage}[b]{\linewidth}\centering
Coefficient of variation
\end{minipage} \\
\midrule\noalign{}
\endhead
\bottomrule\noalign{}
\endlastfoot
LO03 & 3254 & 11.3 & 129.9 \\
LO05 & 4262 & 9.5 & 141.1 \\
LO01 & 4320 & 8.8 & 157.9 \\
LO04 & 5058 & 9.8 & 192.4 \\
LP02 & 17180 & 7.1 & 323.5 \\
LP01 & 18726 & 6.4 & 343.2 \\
LO02 & 9064 & 5.6 & 324.7 \\
\end{longtable}

\section{Conclusion}\label{conclusion}

This in-depth analysis of Campus Lanaken provides a detailed and
data-driven view of how surgical activity evolves from planning to
execution within the operating theatre. The report systematically
examines timing accuracy, room allocation, after-hours activity, urgent
case dynamics, and turnover performance, offering insight into both the
strengths and operational challenges of daily surgical practice.

Across the different sections, the results show that planned and
observed durations are generally well aligned, though variability
between procedures remains considerable. Start-time deviations are
frequent but relatively contained, and their accumulation throughout the
day can still lead to later-than-expected finishes. While the overall
schedule performs reliably, these small deviations highlight the
operational sensitivity of tightly sequenced programs.

After-hours activity is rare at Lanaken, with only about 0.5 percent of
surgeries extending beyond regular hours. Most overruns are brief, yet a
small number of outliers contribute disproportionately to total overtime
minutes. Room assignments remain very stable, with roughly 7.5 percent
of surgeries being relocated. These relocations involve mostly short,
interchangeable procedures and therefore create limited disruption to
the broader workflow.

Urgent cases represent a small fraction of overall activity but still
create meaningful operational interaction with the elective schedule.
Overlaps between urgent and elective cases occur regularly, though their
impact on elective start times is modest and varies across operating
rooms. In most rooms, overlapped elective cases start slightly earlier
or only marginally later than planned, suggesting that urgent
interventions are absorbed efficiently into the daily workflow.

The analysis of shift transitions and idle times further illustrates an
efficient process. Only about 1 percent of surgeries are performed in a
different shift than originally planned, indicating that schedule drift
remains contained. Idle periods between surgeries are very short, with a
median of 0 minutes and an average of 3.2 minutes, confirming
exceptionally fast turnover performance across rooms. A few rooms
display broader idle-time distributions, reflecting more varied case
sequencing or temporary bottlenecks, but these represent exceptions
rather than systemic issues.

Altogether, this document offers a comprehensive quantitative picture of
operating room performance at Campus Lanaken. It shows how surgical
planning translates into real-world execution within a largely stable
and highly efficient environment. The findings confirm strong overall
performance, with targeted opportunities to further stabilise start
times and maintain smooth coordination between elective and urgent
surgical activity.

\end{document}
